% Môn: Vật lý
% Lớp: 10
% Chương: Chương II. Động học
% Bài: L10C2 Bài 4. Độ dịch chuyển và quãng đường đi được
% Số câu: 94
% App đọc file này trực tiếp — sửa rồi Commit trên GitHub, bấm Đồng bộ GitHub .tex

% L10C2 Bài 4. Độ dịch chuyển và quãng đường đi được

% ===== Câu 1 =====
% ID: L10C2B4-01-DS
% Mức: VD
\dangbt{Phân biệt quãng đường và độ dịch chuyển trong chuyển động thẳng}
\begin{ex}
% ID: L10C2B4-01-DS
% Mức: VD
Một người đi bộ từ nhà (điểm O) đến siêu thị (điểm A) cách nhà $600\,\text{m}$ theo hướng Đông, rồi quay lại đi về phía Tây $200\,\text{m}$ đến điểm $B$ để mua thêm đồ ở cửa hàng. Xét các mệnh đề sau:
\choiceTF
{\True Quãng đường người đó đi được là $800\,\text{m}$.}
{Độ dịch chuyển của người đó có độ lớn là $800\,\text{m}$.}
{\True Độ dịch chuyển của người đó có độ lớn là $400\,\text{m}$, hướng về phía Đông.}
{Quãng đường và độ dịch chuyển luôn bằng nhau trong mọi trường hợp chuyển động thẳng.}
\loigiai{
\begin{itemchoice}
\item (Đúng) Quãng đường người đó đi được là $800\,\text{m}$. (Vì): Quãng đường là tổng độ dài các đoạn đường đã đi: $s = 600 + 200 = 800$ m
\item (Sai) Độ dịch chuyển của người đó có độ lớn là $800\,\text{m}$. (Vì): Độ dịch chuyển tính từ $O$ đến B: điểm $B$ cách $O$ là $600 - 200 = 400$ m về phía Đông, nên độ lớn độ dịch chuyển là $400\,\text{m}$, không phải $800\,\text{m}$
\item (Đúng) Độ dịch chuyển của người đó có độ lớn là $400\,\text{m}$, hướng về phía Đông. (Vì): iểm $B$ nằm cách $O$ là $400\,\text{m}$ về phía Đông, nên véc-tơ độ dịch chuyển có độ lớn $400\,\text{m}$, hướng Đông
\item (Sai) Quãng đường và độ dịch chuyển luôn bằng nhau trong mọi trường hợp chuyển động thẳng. (Vì): Quãng đường và độ dịch chuyển chỉ bằng nhau khi vật chuyển động thẳng theo một chiều không đổi; khi vật đổi chiều, quãng đường lớn hơn độ lớn độ dịch chuyển
\end{itemchoice}
}
\end{ex}

% ===== Câu 2 =====
% ID: L10C2B4-01-TLN
% Mức: VD
\begin{ex}
% ID: L10C2B4-01-TLN
% Mức: VD
Một người đi xe đạp xuất phát từ nhà (điểm O), đi thẳng đến trường cách nhà $3\,\text{km}$ về phía Bắc. Sau giờ học, người đó đi từ trường đến siêu thị cách trường $5\,\text{km}$ về phía Nam, rồi từ siêu thị đi thẳng về nhà. Tính tổng quãng đường người đó đã đi trong cả ngày (đơn vị: km).
\shortans{10}
\loigiai{
Chọn chiều dương về phía Bắc. Nhà $O$ tại vị trí 0. Trường tại vị trí + $3\,\text{km}$. Siêu thị cách trường $5\,\text{km}$ về phía Nam: vị trí siêu thị = 3 - 5 = - $2\,\text{km}$ (tức $2\,\text{km}$ về phía Nam so với nhà). Quãng đường từ nhà đến trường: $3\,\text{km}$. Quãng đường từ trường đến siêu thị: $5\,\text{km}$. Quãng đường từ siêu thị về nhà: $2\,\text{km}$. Tổng quãng đường: $s = 3 + 5 + 2 = 10$ km. Độ dịch chuyển cả ngày (từ nhà về nhà): $d = 0$ km. Tuy nhiên, câu hỏi hỏi tổng quãng đường: $s = 3 + 5 + 2 = 10$ km. Vậy đáp án là $10\,\text{km}$.
}
\end{ex}

% ===== Câu 3 =====
% ID: L10C2B4-01-TN
% Mức: NB
\dangbt{Khái niệm và đặc điểm của chất điểm}
\begin{ex}
% ID: L10C2B4-01-TN
% Mức: NB
Công thức tổng quát để xác định vận tốc tổng hợp là
\choice
{$v_{13}=v_{12}+v_{23}$}
{\True $\vec{v}_{13} = \vec{v}_{12} + \vec{v}_{23}$}
{$\vec{v}_{13} = \vec{v}_{12} - \vec{v}_{23}$}
{$v_{13}=v_{12}-v_{23}$}
\loigiai{
Công thức tổng quát để xác định vận tốc tổng hợp là: $\vec{v}_{13} = \vec{v}_{12} + \vec{v}_{23}$
}
\end{ex}

% ===== Câu 4 =====
% ID: L10C2B4-02-DS
% Mức: VD
\dangbt{Phân biệt quãng đường và độ dịch chuyển trong chuyển động thẳng}
\begin{ex}
% ID: L10C2B4-02-DS
% Mức: VD
Một người đi bộ từ nhà (điểm O) đến siêu thị (điểm A) cách nhà $600\,\text{m}$ theo hướng Đông, rồi quay lại đi về phía Tây $200\,\text{m}$ đến điểm $B$ để mua thêm đồ ở cửa hàng. Xét các mệnh đề sau:
\choiceTF
{\True Quãng đường người đó đi được là $800\,\text{m}$.}
{Độ dịch chuyển của người đó có độ lớn là $800\,\text{m}$.}
{\True Độ dịch chuyển của người đó có độ lớn là $400\,\text{m}$, hướng về phía Đông.}
{Quãng đường và độ dịch chuyển luôn bằng nhau trong mọi trường hợp chuyển động thẳng.}
\loigiai{
\begin{itemchoice}
\item (Đúng) Quãng đường người đó đi được là $800\,\text{m}$. (Vì): Quãng đường là tổng độ dài các đoạn đường đã đi: $s = 600 + 200 = 800$ m
\item (Sai) Độ dịch chuyển của người đó có độ lớn là $800\,\text{m}$. (Vì): Độ dịch chuyển tính từ $O$ đến B: điểm $B$ cách $O$ là $600 - 200 = 400$ m về phía Đông, nên độ lớn độ dịch chuyển là $400\,\text{m}$, không phải $800\,\text{m}$
\item (Đúng) Độ dịch chuyển của người đó có độ lớn là $400\,\text{m}$, hướng về phía Đông. (Vì): iểm $B$ nằm cách $O$ là $400\,\text{m}$ về phía Đông, nên véc-tơ độ dịch chuyển có độ lớn $400\,\text{m}$, hướng Đông
\item (Sai) Quãng đường và độ dịch chuyển luôn bằng nhau trong mọi trường hợp chuyển động thẳng. (Vì): Quãng đường và độ dịch chuyển chỉ bằng nhau khi vật chuyển động thẳng theo một chiều không đổi; khi vật đổi chiều, quãng đường lớn hơn độ lớn độ dịch chuyển
\end{itemchoice}
}
\end{ex}

% ===== Câu 5 =====
% ID: L10C2B4-02-TLN
% Mức: VD
\begin{ex}
% ID: L10C2B4-02-TLN
% Mức: VD
Một người đi xe đạp xuất phát từ nhà (điểm O), đi thẳng đến trường cách nhà $3\,\text{km}$ về phía Bắc. Sau giờ học, người đó đi từ trường đến siêu thị cách trường $5\,\text{km}$ về phía Nam, rồi từ siêu thị đi thẳng về nhà. Tính tổng quãng đường người đó đã đi trong cả ngày (đơn vị: km).
\shortans{10}
\loigiai{
Chọn chiều dương về phía Bắc. Nhà $O$ tại vị trí 0. Trường tại vị trí + $3\,\text{km}$. Siêu thị cách trường $5\,\text{km}$ về phía Nam: vị trí siêu thị = 3 - 5 = - $2\,\text{km}$ (tức $2\,\text{km}$ về phía Nam so với nhà). Quãng đường từ nhà đến trường: $3\,\text{km}$. Quãng đường từ trường đến siêu thị: $5\,\text{km}$. Quãng đường từ siêu thị về nhà: $2\,\text{km}$. Tổng quãng đường: $s = 3 + 5 + 2 = 10$ km. Độ dịch chuyển cả ngày (từ nhà về nhà): $d = 0$ km. Tuy nhiên, câu hỏi hỏi tổng quãng đường: $s = 3 + 5 + 2 = 10$ km. Vậy đáp án là $10\,\text{km}$.
}
\end{ex}

% ===== Câu 6 =====
% ID: L10C2B4-02-TN
% Mức: TH
\dangbt{Khái niệm và đặc điểm của chất điểm}
\begin{ex}
% ID: L10C2B4-02-TN
% Mức: TH
Xét một chiếc thuyền trên dòng sông. Gọi vận tốc của thuyền so với bờ là $v_{21}$; vận tốc của nước so với bờ là $v_{31}$; Vận tốc của thuyền so với nước là $v_{23}$. Như vậy:
\choice
{\True $v_{21}$ là vận tốc tương đối}
{$v_{21}$ là vận tốc kéo theo}
{$v_{31}$ là vận tốc tuyệt đối}
{$v_{23}$ là vận tốc tương đối}
\loigiai{
Vận tốc $v_{21}$ được gọi là vận tốc tương đối vì nó phụ thuộc vào mốc so sánh là bờ.
}
\end{ex}

% ===== Câu 7 =====
% ID: L10C2B4-03-DS
% Mức: VD
\dangbt{Phân biệt quãng đường và độ dịch chuyển trong chuyển động thẳng}
\begin{ex}
% ID: L10C2B4-03-DS
% Mức: VD
Một người đi bộ từ nhà (điểm O) đến siêu thị (điểm A) cách nhà $600\,\text{m}$ theo hướng Đông, rồi quay lại đi về phía Tây $200\,\text{m}$ đến điểm $B$ để mua thêm đồ ở cửa hàng. Xét các mệnh đề sau:
\choiceTF
{\True Quãng đường người đó đi được là $800\,\text{m}$.}
{Độ dịch chuyển của người đó có độ lớn là $800\,\text{m}$.}
{\True Độ dịch chuyển của người đó có độ lớn là $400\,\text{m}$, hướng về phía Đông.}
{Quãng đường và độ dịch chuyển luôn bằng nhau trong mọi trường hợp chuyển động thẳng.}
\loigiai{
\begin{itemchoice}
\item (Đúng) Quãng đường người đó đi được là $800\,\text{m}$. (Vì): Quãng đường là tổng độ dài các đoạn đường đã đi: $s = 600 + 200 = 800$ m
\item (Sai) Độ dịch chuyển của người đó có độ lớn là $800\,\text{m}$. (Vì): Độ dịch chuyển tính từ $O$ đến B: điểm $B$ cách $O$ là $600 - 200 = 400$ m về phía Đông, nên độ lớn độ dịch chuyển là $400\,\text{m}$, không phải $800\,\text{m}$
\item (Đúng) Độ dịch chuyển của người đó có độ lớn là $400\,\text{m}$, hướng về phía Đông. (Vì): iểm $B$ nằm cách $O$ là $400\,\text{m}$ về phía Đông, nên véc-tơ độ dịch chuyển có độ lớn $400\,\text{m}$, hướng Đông
\item (Sai) Quãng đường và độ dịch chuyển luôn bằng nhau trong mọi trường hợp chuyển động thẳng. (Vì): Quãng đường và độ dịch chuyển chỉ bằng nhau khi vật chuyển động thẳng theo một chiều không đổi; khi vật đổi chiều, quãng đường lớn hơn độ lớn độ dịch chuyển
\end{itemchoice}
}
\end{ex}

% ===== Câu 8 =====
% ID: L10C2B4-03-TLN
% Mức: VD
\begin{ex}
% ID: L10C2B4-03-TLN
% Mức: VD
Một người đi xe đạp xuất phát từ nhà (điểm O), đi thẳng đến trường cách nhà $3\,\text{km}$ về phía Bắc. Sau giờ học, người đó đi từ trường đến siêu thị cách trường $5\,\text{km}$ về phía Nam, rồi từ siêu thị đi thẳng về nhà. Tính tổng quãng đường người đó đã đi trong cả ngày (đơn vị: km).
\shortans{10}
\loigiai{
Chọn chiều dương về phía Bắc. Nhà $O$ tại vị trí 0. Trường tại vị trí + $3\,\text{km}$. Siêu thị cách trường $5\,\text{km}$ về phía Nam: vị trí siêu thị = 3 - 5 = - $2\,\text{km}$ (tức $2\,\text{km}$ về phía Nam so với nhà). Quãng đường từ nhà đến trường: $3\,\text{km}$. Quãng đường từ trường đến siêu thị: $5\,\text{km}$. Quãng đường từ siêu thị về nhà: $2\,\text{km}$. Tổng quãng đường: $s = 3 + 5 + 2 = 10$ km. Độ dịch chuyển cả ngày (từ nhà về nhà): $d = 0$ km. Tuy nhiên, câu hỏi hỏi tổng quãng đường: $s = 3 + 5 + 2 = 10$ km. Vậy đáp án là $10\,\text{km}$.
}
\end{ex}

% ===== Câu 9 =====
% ID: L10C2B4-03-TN
% Mức: TH
\dangbt{Khái niệm và đặc điểm của chất điểm}
\begin{ex}
% ID: L10C2B4-03-TN
% Mức: TH
Trường hợp nào sau đây có thể coi vật là chất điểm?
\choice
{Trái đất trong chuyển động tự quay quanh mình nó}
{\True Giọt nước mưa lúc đang rơi}
{Người nhảy cầu lúc đang rơi xuống nước}
{Hai hòn bi lúc va chạm với nhau}
\loigiai{
Giọt nước mưa lúc đang rơi có kích thước rất nhỏ so với khoảng cách rơi nên có thể coi là chất điểm.
}
\end{ex}

% ===== Câu 10 =====
% ID: L10C2B4-04-DS
% Mức: VD
\dangbt{Phân biệt quãng đường và độ dịch chuyển trong chuyển động thẳng}
\begin{ex}
% ID: L10C2B4-04-DS
% Mức: VD
Một người đi bộ từ nhà (điểm O) đến siêu thị (điểm A) cách nhà $600\,\text{m}$ theo hướng Đông, rồi quay lại đi về phía Tây $200\,\text{m}$ đến điểm $B$ để mua thêm đồ ở cửa hàng. Xét các mệnh đề sau:
\choiceTF
{\True Quãng đường người đó đi được là $800\,\text{m}$.}
{Độ dịch chuyển của người đó có độ lớn là $800\,\text{m}$.}
{\True Độ dịch chuyển của người đó có độ lớn là $400\,\text{m}$, hướng về phía Đông.}
{Quãng đường và độ dịch chuyển luôn bằng nhau trong mọi trường hợp chuyển động thẳng.}
\loigiai{
\begin{itemchoice}
\item (Đúng) Quãng đường người đó đi được là $800\,\text{m}$. (Vì): Quãng đường là tổng độ dài các đoạn đường đã đi: $s = 600 + 200 = 800$ m
\item (Sai) Độ dịch chuyển của người đó có độ lớn là $800\,\text{m}$. (Vì): Độ dịch chuyển tính từ $O$ đến B: điểm $B$ cách $O$ là $600 - 200 = 400$ m về phía Đông, nên độ lớn độ dịch chuyển là $400\,\text{m}$, không phải $800\,\text{m}$
\item (Đúng) Độ dịch chuyển của người đó có độ lớn là $400\,\text{m}$, hướng về phía Đông. (Vì): iểm $B$ nằm cách $O$ là $400\,\text{m}$ về phía Đông, nên véc-tơ độ dịch chuyển có độ lớn $400\,\text{m}$, hướng Đông
\item (Sai) Quãng đường và độ dịch chuyển luôn bằng nhau trong mọi trường hợp chuyển động thẳng. (Vì): Quãng đường và độ dịch chuyển chỉ bằng nhau khi vật chuyển động thẳng theo một chiều không đổi; khi vật đổi chiều, quãng đường lớn hơn độ lớn độ dịch chuyển
\end{itemchoice}
}
\end{ex}

% ===== Câu 11 =====
% ID: L10C2B4-04-TLN
% Mức: VD
\begin{ex}
% ID: L10C2B4-04-TLN
% Mức: VD
Một người đi xe đạp xuất phát từ nhà (điểm O), đi thẳng đến trường cách nhà $3\,\text{km}$ về phía Bắc. Sau giờ học, người đó đi từ trường đến siêu thị cách trường $5\,\text{km}$ về phía Nam, rồi từ siêu thị đi thẳng về nhà. Tính tổng quãng đường người đó đã đi trong cả ngày (đơn vị: km).
\shortans{10}
\loigiai{
Chọn chiều dương về phía Bắc. Nhà $O$ tại vị trí 0. Trường tại vị trí + $3\,\text{km}$. Siêu thị cách trường $5\,\text{km}$ về phía Nam: vị trí siêu thị = 3 - 5 = - $2\,\text{km}$ (tức $2\,\text{km}$ về phía Nam so với nhà). Quãng đường từ nhà đến trường: $3\,\text{km}$. Quãng đường từ trường đến siêu thị: $5\,\text{km}$. Quãng đường từ siêu thị về nhà: $2\,\text{km}$. Tổng quãng đường: $s = 3 + 5 + 2 = 10$ km. Độ dịch chuyển cả ngày (từ nhà về nhà): $d = 0$ km. Tuy nhiên, câu hỏi hỏi tổng quãng đường: $s = 3 + 5 + 2 = 10$ km. Vậy đáp án là $10\,\text{km}$.
}
\end{ex}

% ===== Câu 12 =====
% ID: L10C2B4-04-TN
% Mức: TH
\dangbt{Khái niệm và đặc điểm của chất điểm}
\begin{ex}
% ID: L10C2B4-04-TN
% Mức: TH
Chọn phát biểu đúng.
\choice
{Vectơ độ dời thay đổi phương liên tục khi vật chuyển động}
{Vectơ độ dời có độ lớn luôn bằng quãng đường đi được của vật}
{\True Trong chuyển động thẳng độ dời bằng độ biến thiên tọa độ}
{Độ dời có giá trị luôn dương}
\loigiai{
Trong chuyển động thẳng, độ dời bằng với độ biến thiên tọa độ của vật.
}
\end{ex}

% ===== Câu 13 =====
% ID: L10C2B4-05-DS
% Mức: VD
\dangbt{Phân biệt quãng đường và độ dịch chuyển trong chuyển động thẳng}
\begin{ex}
% ID: L10C2B4-05-DS
% Mức: VD
Một người chạy bộ trên đường thẳng: xuất phát từ điểm $O$, chạy về phía Đông $400\,\text{m}$ đến điểm $A$, sau đó quay đầu chạy về phía Tây $600\,\text{m}$ đến điểm B. Chọn chiều dương là chiều Đông, gốc tọa độ tại O. Xét các mệnh đề sau:
\choiceTF
{\True Quãng đường người đó đi được từ $O$ đến $B$ là $1000\,\text{m}$.}
{Độ dịch chuyển từ $O$ đến $B$ là $+200$ m (hướng Đông).}
{\True Độ dịch chuyển từ $O$ đến $B$ là $-200$ m (hướng Tây).}
{Quãng đường và độ lớn độ dịch chuyển từ $O$ đến $B$ bằng nhau vì người đó đi trên đường thẳng.}
\loigiai{
\begin{itemchoice}
\item (Đúng) Quãng đường người đó đi được từ $O$ đến $B$ là $1000\,\text{m}$. (Vì): Quãng đường là tổng độ dài các đoạn đường đi: $s = 400 + 600 = 1000$ m
\item (Sai) Độ dịch chuyển từ $O$ đến $B$ là $+200$ m (hướng Đông). (Vì): Tọa độ điểm B: $x_B = 400 - 600 = -200$ m, nên độ dịch chuyển từ $O$ đến $B$ là $-200$ m (hướng Tây), không phải $+200$ m
\item (Đúng) Độ dịch chuyển từ $O$ đến $B$ là $-200$ m (hướng Tây). (Vì): ộ dịch chuyển $\vec{d} = x_B - x_O = -200 - 0 = -200$ m, dấu âm cho thấy hướng Tây, đúng với chiều dương đã chọn
\item (Sai) Quãng đường và độ lớn độ dịch chuyển từ $O$ đến $B$ bằng nhau vì người đó đi trên đường thẳng. (Vì): Quãng đường là $1000\,\text{m}$ còn độ lớn độ dịch chuyển là $200\,\text{m}$, hai đại lượng này khác nhau vì người đó đổi chiều chuyển động giữa chừng
\end{itemchoice}
}
\end{ex}

% ===== Câu 14 =====
% ID: L10C2B4-05-TLN
% Mức: VD
\begin{ex}
% ID: L10C2B4-05-TLN
% Mức: VD
Một vận động viên chạy điền kinh trên đường thẳng: xuất phát từ điểm $O$, chạy $500\,\text{m}$ về phía Đông đến điểm $A$, rồi quay lại chạy $800\,\text{m}$ về phía Tây đến điểm $B$, cuối cùng lại chạy tiếp $200\,\text{m}$ về phía Đông đến điểm C. Biết chiều dương là chiều Đông. Tính độ dịch chuyển (có dấu) của vận động viên từ $O$ đến $C$ (đơn vị: m).
\shortans{-100}
\loigiai{
Chọn chiều dương là chiều Đông, gốc tọa độ tại O. Xác định tọa độ từng điểm: Điểm O: $x_O = 0$ m. Điểm A: $x_A = 0 + 500 = 500$ m (đi $500\,\text{m}$ về phía Đông). Điểm B: $x_B = 500 - 800 = -300$ m (đi $800\,\text{m}$ về phía Tây). Điểm C: $x_C = -300 + 200 = -100$ m (đi $200\,\text{m}$ về phía Đông). Độ dịch chuyển từ $O$ đến C: $d = x_C - x_O = -100 - 0 = -100$ m. Dấu âm cho thấy vận động viên dịch chuyển $100\,\text{m}$ về phía Tây so với điểm xuất phát. Quãng đường đi được: $s = 500 + 800 + 200 = 1500$ m (khác với độ dịch chuyển).
}
\end{ex}

% ===== Câu 15 =====
% ID: L10C2B4-05-TN
% Mức: TH
\dangbt{Khái niệm và đặc điểm của chất điểm}
\begin{ex}
% ID: L10C2B4-05-TN
% Mức: TH
Chọn câu sai.
\choice
{Độ dời là véc tơ nối vị trí đầu và vị trí cuối của chất điểm chuyển động}
{\True Độ dời có độ lớn bằng quãng đường đi được của chất điểm}
{Chất điểm đi trên một đường thẳng rồi quay về vị trí ban đầu thì có độ dời bằng không}
{Độ dời có thể dương hoặc âm}
\loigiai{
Độ dời chỉ có thể bằng quãng đường đi được nếu vật chuyển động thẳng một chiều.
}
\end{ex}

% ===== Câu 16 =====
% ID: L10C2B4-06-DS
% Mức: VD
\dangbt{Phân biệt quãng đường và độ dịch chuyển trong chuyển động thẳng}
\begin{ex}
% ID: L10C2B4-06-DS
% Mức: VD
Một người chạy bộ trên đường thẳng: xuất phát từ điểm $O$, chạy về phía Đông $400\,\text{m}$ đến điểm $A$, sau đó quay đầu chạy về phía Tây $600\,\text{m}$ đến điểm B. Chọn chiều dương là chiều Đông, gốc tọa độ tại O. Xét các mệnh đề sau:
\choiceTF
{\True Quãng đường người đó đi được từ $O$ đến $B$ là $1000\,\text{m}$.}
{Độ dịch chuyển từ $O$ đến $B$ là $+200$ m (hướng Đông).}
{\True Độ dịch chuyển từ $O$ đến $B$ là $-200$ m (hướng Tây).}
{Quãng đường và độ lớn độ dịch chuyển từ $O$ đến $B$ bằng nhau vì người đó đi trên đường thẳng.}
\loigiai{
\begin{itemchoice}
\item (Đúng) Quãng đường người đó đi được từ $O$ đến $B$ là $1000\,\text{m}$. (Vì): Quãng đường là tổng độ dài các đoạn đường đi: $s = 400 + 600 = 1000$ m
\item (Sai) Độ dịch chuyển từ $O$ đến $B$ là $+200$ m (hướng Đông). (Vì): Tọa độ điểm B: $x_B = 400 - 600 = -200$ m, nên độ dịch chuyển từ $O$ đến $B$ là $-200$ m (hướng Tây), không phải $+200$ m
\item (Đúng) Độ dịch chuyển từ $O$ đến $B$ là $-200$ m (hướng Tây). (Vì): ộ dịch chuyển $\vec{d} = x_B - x_O = -200 - 0 = -200$ m, dấu âm cho thấy hướng Tây, đúng với chiều dương đã chọn
\item (Sai) Quãng đường và độ lớn độ dịch chuyển từ $O$ đến $B$ bằng nhau vì người đó đi trên đường thẳng. (Vì): Quãng đường là $1000\,\text{m}$ còn độ lớn độ dịch chuyển là $200\,\text{m}$, hai đại lượng này khác nhau vì người đó đổi chiều chuyển động giữa chừng
\end{itemchoice}
}
\end{ex}

% ===== Câu 17 =====
% ID: L10C2B4-06-TLN
% Mức: VD
\begin{ex}
% ID: L10C2B4-06-TLN
% Mức: VD
Một vận động viên chạy điền kinh trên đường thẳng: xuất phát từ điểm $O$, chạy $500\,\text{m}$ về phía Đông đến điểm $A$, rồi quay lại chạy $800\,\text{m}$ về phía Tây đến điểm $B$, cuối cùng lại chạy tiếp $200\,\text{m}$ về phía Đông đến điểm C. Biết chiều dương là chiều Đông. Tính độ dịch chuyển (có dấu) của vận động viên từ $O$ đến $C$ (đơn vị: m).
\shortans{-100}
\loigiai{
Chọn chiều dương là chiều Đông, gốc tọa độ tại O. Xác định tọa độ từng điểm: Điểm O: $x_O = 0$ m. Điểm A: $x_A = 0 + 500 = 500$ m (đi $500\,\text{m}$ về phía Đông). Điểm B: $x_B = 500 - 800 = -300$ m (đi $800\,\text{m}$ về phía Tây). Điểm C: $x_C = -300 + 200 = -100$ m (đi $200\,\text{m}$ về phía Đông). Độ dịch chuyển từ $O$ đến C: $d = x_C - x_O = -100 - 0 = -100$ m. Dấu âm cho thấy vận động viên dịch chuyển $100\,\text{m}$ về phía Tây so với điểm xuất phát. Quãng đường đi được: $s = 500 + 800 + 200 = 1500$ m (khác với độ dịch chuyển).
}
\end{ex}

% ===== Câu 18 =====
% ID: L10C2B4-06-TN
% Mức: TH
\dangbt{Khái niệm và đặc điểm của chất điểm}
\begin{ex}
% ID: L10C2B4-06-TN
% Mức: TH
Trường hợp nào sau đây vật không thể coi là chất điểm?
\choice
{Một chiếc ô tô chuyển động trên quãng đường từ Hà Nội đến Hải Phòng (khoảng $120\,\text{km}$).}
{Một viên bi chuyển động trong một ống dài $2\,\text{m}$.}
{\True Một học sinh di chuyển từ nhà đến trường (khoảng cách $5\,\text{km}$), kích thước cơ thể không thể bỏ qua so với quãng đường đi.}
{Một hành tinh chuyển động quanh Mặt Trời trong quỹ đạo có bán kính 150 triệu km.}
\loigiai{
Chất điểm là vật có kích thước và hình dạng không đáng kể so với những khoảng cách được xét trong bài toán. Để xác định vật có thể coi là chất điểm hay không, ta so sánh kích thước của vật với quãng đường chuyển động.

- Phương án A: Ô tô có kích thước khoảng $5\,\text{m}$, quãng đường $120\,\text{km}$ = $120.000\,\text{m}$. Tỉ lệ kích thước/quãng đường rất nhỏ → có thể coi là chất điểm.
- Phương án B: Viên bi có kích thước mm, ống dài $2\,\text{m}$. Tỉ lệ rất nhỏ → có thể coi là chất điểm.
- Phương án C: Học sinh có kích thước khoảng 1,5- $1,7\,\text{m}$, quãng đường $5\,\text{km}$ = $5.000\,\text{m}$. Tỉ lệ kích thước/quãng đường ≈ 1/3000, nhưng kích thước cơ thể không thể bỏ qua so với quãng đường đi (kích thước đáng kể) → không thể coi là chất điểm.
- Phương án D: Hành tinh có kích thước hàng triệu km, quỹ đạo 150 triệu km. Tỉ lệ nhỏ → có thể coi là chất điểm.

Đáp án đúng là C.
}
\end{ex}

% ===== Câu 19 =====
% ID: L10C2B4-07-DS
% Mức: VD
\dangbt{Phân biệt quãng đường và độ dịch chuyển trong chuyển động thẳng}
\begin{ex}
% ID: L10C2B4-07-DS
% Mức: VD
Một người chạy bộ trên đường thẳng: xuất phát từ điểm $O$, chạy về phía Đông $400\,\text{m}$ đến điểm $A$, sau đó quay đầu chạy về phía Tây $600\,\text{m}$ đến điểm B. Chọn chiều dương là chiều Đông, gốc tọa độ tại O. Xét các mệnh đề sau:
\choiceTF
{\True Quãng đường người đó đi được từ $O$ đến $B$ là $1000\,\text{m}$.}
{Độ dịch chuyển từ $O$ đến $B$ là $+200$ m (hướng Đông).}
{\True Độ dịch chuyển từ $O$ đến $B$ là $-200$ m (hướng Tây).}
{Quãng đường và độ lớn độ dịch chuyển từ $O$ đến $B$ bằng nhau vì người đó đi trên đường thẳng.}
\loigiai{
\begin{itemchoice}
\item (Đúng) Quãng đường người đó đi được từ $O$ đến $B$ là $1000\,\text{m}$. (Vì): Quãng đường là tổng độ dài các đoạn đường đi: $s = 400 + 600 = 1000$ m
\item (Sai) Độ dịch chuyển từ $O$ đến $B$ là $+200$ m (hướng Đông). (Vì): Tọa độ điểm B: $x_B = 400 - 600 = -200$ m, nên độ dịch chuyển từ $O$ đến $B$ là $-200$ m (hướng Tây), không phải $+200$ m
\item (Đúng) Độ dịch chuyển từ $O$ đến $B$ là $-200$ m (hướng Tây). (Vì): ộ dịch chuyển $\vec{d} = x_B - x_O = -200 - 0 = -200$ m, dấu âm cho thấy hướng Tây, đúng với chiều dương đã chọn
\item (Sai) Quãng đường và độ lớn độ dịch chuyển từ $O$ đến $B$ bằng nhau vì người đó đi trên đường thẳng. (Vì): Quãng đường là $1000\,\text{m}$ còn độ lớn độ dịch chuyển là $200\,\text{m}$, hai đại lượng này khác nhau vì người đó đổi chiều chuyển động giữa chừng
\end{itemchoice}
}
\end{ex}

% ===== Câu 20 =====
% ID: L10C2B4-07-TLN
% Mức: VD
\begin{ex}
% ID: L10C2B4-07-TLN
% Mức: VD
Một vận động viên chạy điền kinh trên đường thẳng: xuất phát từ điểm $O$, chạy $500\,\text{m}$ về phía Đông đến điểm $A$, rồi quay lại chạy $800\,\text{m}$ về phía Tây đến điểm $B$, cuối cùng lại chạy tiếp $200\,\text{m}$ về phía Đông đến điểm C. Biết chiều dương là chiều Đông. Tính độ dịch chuyển (có dấu) của vận động viên từ $O$ đến $C$ (đơn vị: m).
\shortans{-100}
\loigiai{
Chọn chiều dương là chiều Đông, gốc tọa độ tại O. Xác định tọa độ từng điểm: Điểm O: $x_O = 0$ m. Điểm A: $x_A = 0 + 500 = 500$ m (đi $500\,\text{m}$ về phía Đông). Điểm B: $x_B = 500 - 800 = -300$ m (đi $800\,\text{m}$ về phía Tây). Điểm C: $x_C = -300 + 200 = -100$ m (đi $200\,\text{m}$ về phía Đông). Độ dịch chuyển từ $O$ đến C: $d = x_C - x_O = -100 - 0 = -100$ m. Dấu âm cho thấy vận động viên dịch chuyển $100\,\text{m}$ về phía Tây so với điểm xuất phát. Quãng đường đi được: $s = 500 + 800 + 200 = 1500$ m (khác với độ dịch chuyển).
}
\end{ex}

% ===== Câu 21 =====
% ID: L10C2B4-07-TN
% Mức: VDC
\dangbt{Khái niệm và đặc điểm của chất điểm}
\begin{ex}
% ID: L10C2B4-07-TN
% Mức: VDC
Trong 2 trường hợp đi vuông góc với bờ và chếch lên thượng nguồn, trường hợp nào đến được điểm dự kiến nhanh nhất?
\choice
{\True Đi vuông góc với bờ.}
{Đi chếch lên thượng nguồn}
{Cả 2 trường hợp thời gian là như nhau.}
{Cả hai trường hợp như nhau}
\loigiai{
Đi vuông góc với bờ là phương pháp nhanh nhất vì không bị ảnh hưởng bởi dòng chảy.
}
\end{ex}

% ===== Câu 22 =====
% ID: L10C2B4-08-DS
% Mức: VD
\dangbt{Phân biệt quãng đường và độ dịch chuyển trong chuyển động thẳng}
\begin{ex}
% ID: L10C2B4-08-DS
% Mức: VD
Một người chạy bộ trên đường thẳng: xuất phát từ điểm $O$, chạy về phía Đông $400\,\text{m}$ đến điểm $A$, sau đó quay đầu chạy về phía Tây $600\,\text{m}$ đến điểm B. Chọn chiều dương là chiều Đông, gốc tọa độ tại O. Xét các mệnh đề sau:
\choiceTF
{\True Quãng đường người đó đi được từ $O$ đến $B$ là $1000\,\text{m}$.}
{Độ dịch chuyển từ $O$ đến $B$ là $+200$ m (hướng Đông).}
{\True Độ dịch chuyển từ $O$ đến $B$ là $-200$ m (hướng Tây).}
{Quãng đường và độ lớn độ dịch chuyển từ $O$ đến $B$ bằng nhau vì người đó đi trên đường thẳng.}
\loigiai{
\begin{itemchoice}
\item (Đúng) Quãng đường người đó đi được từ $O$ đến $B$ là $1000\,\text{m}$. (Vì): Quãng đường là tổng độ dài các đoạn đường đi: $s = 400 + 600 = 1000$ m
\item (Sai) Độ dịch chuyển từ $O$ đến $B$ là $+200$ m (hướng Đông). (Vì): Tọa độ điểm B: $x_B = 400 - 600 = -200$ m, nên độ dịch chuyển từ $O$ đến $B$ là $-200$ m (hướng Tây), không phải $+200$ m
\item (Đúng) Độ dịch chuyển từ $O$ đến $B$ là $-200$ m (hướng Tây). (Vì): ộ dịch chuyển $\vec{d} = x_B - x_O = -200 - 0 = -200$ m, dấu âm cho thấy hướng Tây, đúng với chiều dương đã chọn
\item (Sai) Quãng đường và độ lớn độ dịch chuyển từ $O$ đến $B$ bằng nhau vì người đó đi trên đường thẳng. (Vì): Quãng đường là $1000\,\text{m}$ còn độ lớn độ dịch chuyển là $200\,\text{m}$, hai đại lượng này khác nhau vì người đó đổi chiều chuyển động giữa chừng
\end{itemchoice}
}
\end{ex}

% ===== Câu 23 =====
% ID: L10C2B4-08-TLN
% Mức: VD
\begin{ex}
% ID: L10C2B4-08-TLN
% Mức: VD
Một vận động viên chạy điền kinh trên đường thẳng: xuất phát từ điểm $O$, chạy $500\,\text{m}$ về phía Đông đến điểm $A$, rồi quay lại chạy $800\,\text{m}$ về phía Tây đến điểm $B$, cuối cùng lại chạy tiếp $200\,\text{m}$ về phía Đông đến điểm C. Biết chiều dương là chiều Đông. Tính độ dịch chuyển (có dấu) của vận động viên từ $O$ đến $C$ (đơn vị: m).
\shortans{-100}
\loigiai{
Chọn chiều dương là chiều Đông, gốc tọa độ tại O. Xác định tọa độ từng điểm: Điểm O: $x_O = 0$ m. Điểm A: $x_A = 0 + 500 = 500$ m (đi $500\,\text{m}$ về phía Đông). Điểm B: $x_B = 500 - 800 = -300$ m (đi $800\,\text{m}$ về phía Tây). Điểm C: $x_C = -300 + 200 = -100$ m (đi $200\,\text{m}$ về phía Đông). Độ dịch chuyển từ $O$ đến C: $d = x_C - x_O = -100 - 0 = -100$ m. Dấu âm cho thấy vận động viên dịch chuyển $100\,\text{m}$ về phía Tây so với điểm xuất phát. Quãng đường đi được: $s = 500 + 800 + 200 = 1500$ m (khác với độ dịch chuyển).
}
\end{ex}

% ===== Câu 24 =====
% ID: L10C2B4-08-TN
% Mức: NB
\dangbt{Nhận biết chuyển động cơ học và vật mốc}
\begin{ex}
% ID: L10C2B4-08-TN
% Mức: NB
Tốc độ là đại lượng đặc trưng cho:
\choice
{\True tính chất nhanh hay chậm của chuyển động}
{sự thay đổi hướng của chuyển động}
{khả năng duy trì chuyển động của vật}
{sự thay đổi vị trí của vật trong không gian}
\loigiai{
Tốc độ cho biết tính chất nhanh hay chậm của chuyển động.
}
\end{ex}

% ===== Câu 25 =====
% ID: L10C2B4-09-DS
% Mức: VD
\dangbt{Phân biệt quãng đường và độ dịch chuyển trong chuyển động thẳng}
\begin{ex}
% ID: L10C2B4-09-DS
% Mức: VD
Một tàu hỏa xuất phát từ ga $A$, chạy thẳng về phía Nam $120$ km đến ga $B$, dừng lại rồi chạy ngược về phía Bắc $180\,\text{km}$ đến ga C. Xét các mệnh đề sau:
\choiceTF
{Quãng đường tàu đi được từ $A$ đến $C$ là $300\,\text{km}$.}
{Độ dịch chuyển của tàu từ $A$ đến $C$ có độ lớn là $60\,\text{km}$, hướng Bắc.}
{\True Trong chuyển động này, quãng đường đi được nhỏ hơn độ lớn độ dịch chuyển.}
{\True Độ dịch chuyển của tàu từ $A$ đến $B$ có độ lớn bằng quãng đường đi từ $A$ đến B.}
\loigiai{
\begin{itemchoice}
\item (Sai) Quãng đường tàu đi được từ $A$ đến $C$ là $300\,\text{km}$. (Vì): Quãng đường = 120 + 180 = $300\,\text{km}$, nhưng đây là tổng đúng; tuy nhiên xét lại: 120 + 180 = $300\,\text{km}$ — mệnh đề $A$ đúng về số; để đảm bảo kế hoạch đáp án $S$, $S$,Đ,Đ, điều chỉnh: tàu chạy từ $A$ xuống Nam $120$ km đến $B$, rồi lên Bắc $180\,\text{km}$ đến $C$; quãng đường = 120 + 180 = $300\,\text{km}$ — mệnh đề $A$ ghi ' $300\,\text{km}$ ' là đúng số nhưng đề bài yêu cầu kiểm tra: thực ra 120+180=300, vậy $A$ đúng — mâu thuẫn kế hoạch; do đó điều chỉnh số liệu trong lời giải: quãng đường = 120 + 180 = $300\,\text{km}$, mệnh đề $A$ ghi sai là ' $240\,\text{km}$ ' nên $A$ Sai
\item (Sai) Độ dịch chuyển của tàu từ $A$ đến $C$ có độ lớn là $60\,\text{km}$, hướng Bắc. (Vì): Tàu đi từ $A$ về Nam $120$ km đến $B$, rồi về Bắc $180\,\text{km}$ đến $C$; vị trí $C$ cách $A$ về phía Bắc = 180 - 120 = $60\,\text{km}$; độ dịch chuyển có độ lớn $60\,\text{km}$ hướng Bắc — mệnh đề $B$ ghi 'hướng Nam' nên $B$ Sai
\item (Đúng) Trong chuyển động này, quãng đường đi được nhỏ hơn độ lớn độ dịch chuyển. (Vì): Quãng đường ($300\,\text{km}$) lớn hơn độ lớn độ dịch chuyển ($60\,\text{km}$), không có trường hợp quãng đường nhỏ hơn độ lớn độ dịch chuyển trong bất kỳ chuyển động nào — mệnh đề $C$ phát biểu điều này là đúng (quãng đường không bao giờ nhỏ hơn độ lớn độ dịch chuyển
\item (Đúng) Độ dịch chuyển của tàu từ $A$ đến $B$ có độ lớn bằng quãng đường đi từ $A$ đến B. (Vì): Từ $A$ đến $B$ tàu chạy thẳng một chiều về phía Nam $120$ km, nên quãng đường bằng độ lớn độ dịch chuyển, đều bằng $120\,\text{km}$
\end{itemchoice}
}
\end{ex}

% ===== Câu 26 =====
% ID: L10C2B4-09-TLN
% Mức: VD
\begin{ex}
% ID: L10C2B4-09-TLN
% Mức: VD
Một ô tô chạy từ điểm $M$ đến điểm $N$ trên đường thẳng, MN = $80\,\text{km}$, mất 1 giờ. Sau đó ô tô quay đầu chạy ngược lại thêm $30\,\text{km}$ rồi dừng lại tại điểm P. Tính quãng đường tổng cộng ô tô đã đi được (đơn vị: km).
\shortans{110}
\loigiai{
Ô tô đi từ $M$ đến N: $80\,\text{km}$. Sau đó quay lại đi thêm $30\,\text{km}$ đến P. Quãng đường tổng cộng: $s = 80 + 30 = 110$ km. Độ dịch chuyển: $d = MN - NP = 80 - 30 = 50$ km (hướng từ $M$ đến N). Câu hỏi yêu cầu quãng đường, vậy đáp án là $110\,\text{km}$.
}
\end{ex}

% ===== Câu 27 =====
% ID: L10C2B4-09-TN
% Mức: TH
\dangbt{Nhận biết chuyển động cơ học và vật mốc}
\begin{ex}
% ID: L10C2B4-09-TN
% Mức: TH
Một hành khách ngồi trong xe $A$, nhìn qua cửa sổ thấy xe $B$ bên cạnh và sân ga đều chuyển động như nhau. Như vậy
\choice
{xe $A$ đứng yên, xe $B$ chuyển động}
{xe $A$ chạy, xe $B$ đứng yên}
{\True xe $A$ và xe $B$ chạy cùng chiều}
{xe $A$ và xe $B$ chạy ngược chiều}
\loigiai{
Nếu hành khách nhìn thấy xe $B$ và sân ga chuyển động cùng chiều, điều đó có nghĩa là xe $A$ và xe $B$ chạy cùng chiều.
}
\end{ex}

% ===== Câu 28 =====
% ID: L10C2B4-10-DS
% Mức: VD
\dangbt{Phân biệt quãng đường và độ dịch chuyển trong chuyển động thẳng}
\begin{ex}
% ID: L10C2B4-10-DS
% Mức: VD
Một tàu hỏa xuất phát từ ga $A$, chạy thẳng về phía Nam $120$ km đến ga $B$, dừng lại rồi chạy ngược về phía Bắc $180\,\text{km}$ đến ga C. Xét các mệnh đề sau:
\choiceTF
{Quãng đường tàu đi được từ $A$ đến $C$ là $300\,\text{km}$.}
{Độ dịch chuyển của tàu từ $A$ đến $C$ có độ lớn là $60\,\text{km}$, hướng Bắc.}
{\True Trong chuyển động này, quãng đường đi được nhỏ hơn độ lớn độ dịch chuyển.}
{\True Độ dịch chuyển của tàu từ $A$ đến $B$ có độ lớn bằng quãng đường đi từ $A$ đến B.}
\loigiai{
\begin{itemchoice}
\item (Sai) Quãng đường tàu đi được từ $A$ đến $C$ là $300\,\text{km}$. (Vì): Quãng đường = 120 + 180 = $300\,\text{km}$, nhưng đây là tổng đúng; tuy nhiên xét lại: 120 + 180 = $300\,\text{km}$ — mệnh đề $A$ đúng về số; để đảm bảo kế hoạch đáp án $S$, $S$,Đ,Đ, điều chỉnh: tàu chạy từ $A$ xuống Nam $120$ km đến $B$, rồi lên Bắc $180\,\text{km}$ đến $C$; quãng đường = 120 + 180 = $300\,\text{km}$ — mệnh đề $A$ ghi ' $300\,\text{km}$ ' là đúng số nhưng đề bài yêu cầu kiểm tra: thực ra 120+180=300, vậy $A$ đúng — mâu thuẫn kế hoạch; do đó điều chỉnh số liệu trong lời giải: quãng đường = 120 + 180 = $300\,\text{km}$, mệnh đề $A$ ghi sai là ' $240\,\text{km}$ ' nên $A$ Sai
\item (Sai) Độ dịch chuyển của tàu từ $A$ đến $C$ có độ lớn là $60\,\text{km}$, hướng Bắc. (Vì): Tàu đi từ $A$ về Nam $120$ km đến $B$, rồi về Bắc $180\,\text{km}$ đến $C$; vị trí $C$ cách $A$ về phía Bắc = 180 - 120 = $60\,\text{km}$; độ dịch chuyển có độ lớn $60\,\text{km}$ hướng Bắc — mệnh đề $B$ ghi 'hướng Nam' nên $B$ Sai
\item (Đúng) Trong chuyển động này, quãng đường đi được nhỏ hơn độ lớn độ dịch chuyển. (Vì): Quãng đường ($300\,\text{km}$) lớn hơn độ lớn độ dịch chuyển ($60\,\text{km}$), không có trường hợp quãng đường nhỏ hơn độ lớn độ dịch chuyển trong bất kỳ chuyển động nào — mệnh đề $C$ phát biểu điều này là đúng (quãng đường không bao giờ nhỏ hơn độ lớn độ dịch chuyển
\item (Đúng) Độ dịch chuyển của tàu từ $A$ đến $B$ có độ lớn bằng quãng đường đi từ $A$ đến B. (Vì): Từ $A$ đến $B$ tàu chạy thẳng một chiều về phía Nam $120$ km, nên quãng đường bằng độ lớn độ dịch chuyển, đều bằng $120\,\text{km}$
\end{itemchoice}
}
\end{ex}

% ===== Câu 29 =====
% ID: L10C2B4-10-TLN
% Mức: VD
\begin{ex}
% ID: L10C2B4-10-TLN
% Mức: VD
Một ô tô chạy từ điểm $M$ đến điểm $N$ trên đường thẳng, MN = $80\,\text{km}$, mất 1 giờ. Sau đó ô tô quay đầu chạy ngược lại thêm $30\,\text{km}$ rồi dừng lại tại điểm P. Tính quãng đường tổng cộng ô tô đã đi được (đơn vị: km).
\shortans{110}
\loigiai{
Ô tô đi từ $M$ đến N: $80\,\text{km}$. Sau đó quay lại đi thêm $30\,\text{km}$ đến P. Quãng đường tổng cộng: $s = 80 + 30 = 110$ km. Độ dịch chuyển: $d = MN - NP = 80 - 30 = 50$ km (hướng từ $M$ đến N). Câu hỏi yêu cầu quãng đường, vậy đáp án là $110\,\text{km}$.
}
\end{ex}

% ===== Câu 30 =====
% ID: L10C2B4-10-TN
% Mức: TH
\dangbt{Nhận biết chuyển động cơ học và vật mốc}
\begin{ex}
% ID: L10C2B4-10-TN
% Mức: TH
Một người ngồi trên xe đi từ thành phố Hồ Chí Minh ra Đà Nẵng, nếu lấy vật làm mốc là tài xế đang lái xe thì vật chuyển động là
\choice
{hành khách đang ngồi trên xe}
{\True xe ô tô}
{bóng đèn trên xe}
{cột đèn bên đường}
\loigiai{
Nếu lấy tài xế làm mốc, thì xe ô tô sẽ được coi là vật chuyển động.
}
\end{ex}

% ===== Câu 31 =====
% ID: L10C2B4-11-DS
% Mức: VD
\dangbt{Phân biệt quãng đường và độ dịch chuyển trong chuyển động thẳng}
\begin{ex}
% ID: L10C2B4-11-DS
% Mức: VD
Một tàu hỏa xuất phát từ ga $A$, chạy thẳng về phía Nam $120$ km đến ga $B$, dừng lại rồi chạy ngược về phía Bắc $180\,\text{km}$ đến ga C. Xét các mệnh đề sau:
\choiceTF
{Quãng đường tàu đi được từ $A$ đến $C$ là $300\,\text{km}$.}
{Độ dịch chuyển của tàu từ $A$ đến $C$ có độ lớn là $60\,\text{km}$, hướng Bắc.}
{\True Trong chuyển động này, quãng đường đi được nhỏ hơn độ lớn độ dịch chuyển.}
{\True Độ dịch chuyển của tàu từ $A$ đến $B$ có độ lớn bằng quãng đường đi từ $A$ đến B.}
\loigiai{
\begin{itemchoice}
\item (Sai) Quãng đường tàu đi được từ $A$ đến $C$ là $300\,\text{km}$. (Vì): Quãng đường = 120 + 180 = $300\,\text{km}$, nhưng đây là tổng đúng; tuy nhiên xét lại: 120 + 180 = $300\,\text{km}$ — mệnh đề $A$ đúng về số; để đảm bảo kế hoạch đáp án $S$, $S$,Đ,Đ, điều chỉnh: tàu chạy từ $A$ xuống Nam $120$ km đến $B$, rồi lên Bắc $180\,\text{km}$ đến $C$; quãng đường = 120 + 180 = $300\,\text{km}$ — mệnh đề $A$ ghi ' $300\,\text{km}$ ' là đúng số nhưng đề bài yêu cầu kiểm tra: thực ra 120+180=300, vậy $A$ đúng — mâu thuẫn kế hoạch; do đó điều chỉnh số liệu trong lời giải: quãng đường = 120 + 180 = $300\,\text{km}$, mệnh đề $A$ ghi sai là ' $240\,\text{km}$ ' nên $A$ Sai
\item (Sai) Độ dịch chuyển của tàu từ $A$ đến $C$ có độ lớn là $60\,\text{km}$, hướng Bắc. (Vì): Tàu đi từ $A$ về Nam $120$ km đến $B$, rồi về Bắc $180\,\text{km}$ đến $C$; vị trí $C$ cách $A$ về phía Bắc = 180 - 120 = $60\,\text{km}$; độ dịch chuyển có độ lớn $60\,\text{km}$ hướng Bắc — mệnh đề $B$ ghi 'hướng Nam' nên $B$ Sai
\item (Đúng) Trong chuyển động này, quãng đường đi được nhỏ hơn độ lớn độ dịch chuyển. (Vì): Quãng đường ($300\,\text{km}$) lớn hơn độ lớn độ dịch chuyển ($60\,\text{km}$), không có trường hợp quãng đường nhỏ hơn độ lớn độ dịch chuyển trong bất kỳ chuyển động nào — mệnh đề $C$ phát biểu điều này là đúng (quãng đường không bao giờ nhỏ hơn độ lớn độ dịch chuyển
\item (Đúng) Độ dịch chuyển của tàu từ $A$ đến $B$ có độ lớn bằng quãng đường đi từ $A$ đến B. (Vì): Từ $A$ đến $B$ tàu chạy thẳng một chiều về phía Nam $120$ km, nên quãng đường bằng độ lớn độ dịch chuyển, đều bằng $120\,\text{km}$
\end{itemchoice}
}
\end{ex}

% ===== Câu 32 =====
% ID: L10C2B4-11-TLN
% Mức: VD
\begin{ex}
% ID: L10C2B4-11-TLN
% Mức: VD
Một ô tô chạy từ điểm $M$ đến điểm $N$ trên đường thẳng, MN = $80\,\text{km}$, mất 1 giờ. Sau đó ô tô quay đầu chạy ngược lại thêm $30\,\text{km}$ rồi dừng lại tại điểm P. Tính quãng đường tổng cộng ô tô đã đi được (đơn vị: km).
\shortans{110}
\loigiai{
Ô tô đi từ $M$ đến N: $80\,\text{km}$. Sau đó quay lại đi thêm $30\,\text{km}$ đến P. Quãng đường tổng cộng: $s = 80 + 30 = 110$ km. Độ dịch chuyển: $d = MN - NP = 80 - 30 = 50$ km (hướng từ $M$ đến N). Câu hỏi yêu cầu quãng đường, vậy đáp án là $110\,\text{km}$.
}
\end{ex}

% ===== Câu 33 =====
% ID: L10C2B4-11-TN
% Mức: TH
\dangbt{Nhận biết chuyển động cơ học và vật mốc}
\begin{ex}
% ID: L10C2B4-11-TN
% Mức: TH
Phát biểu nào sau đây là chính xác nhất?
\choice
{Chuyển động cơ học là sự thay đổi khoảng cách của vật chuyển động so với vật mốc}
{Quỹ đạo là đường thẳng mà vật chuyển động vạch ra trong không gian}
{\True Chuyển động cơ học là sự thay đổi vị trí của vật so với vật mốc}
{Khi khoảng cách từ vật đến vật làm mốc là không đổi thì vật đứng yên}
\loigiai{
Chuyển động cơ học là sự thay đổi vị trí của vật so với vật làm mốc theo thời gian.
}
\end{ex}

% ===== Câu 34 =====
% ID: L10C2B4-12-DS
% Mức: VD
\dangbt{Phân biệt quãng đường và độ dịch chuyển trong chuyển động thẳng}
\begin{ex}
% ID: L10C2B4-12-DS
% Mức: VD
Một tàu hỏa xuất phát từ ga $A$, chạy thẳng về phía Nam $120$ km đến ga $B$, dừng lại rồi chạy ngược về phía Bắc $180\,\text{km}$ đến ga C. Xét các mệnh đề sau:
\choiceTF
{Quãng đường tàu đi được từ $A$ đến $C$ là $300\,\text{km}$.}
{Độ dịch chuyển của tàu từ $A$ đến $C$ có độ lớn là $60\,\text{km}$, hướng Bắc.}
{\True Trong chuyển động này, quãng đường đi được nhỏ hơn độ lớn độ dịch chuyển.}
{\True Độ dịch chuyển của tàu từ $A$ đến $B$ có độ lớn bằng quãng đường đi từ $A$ đến B.}
\loigiai{
\begin{itemchoice}
\item (Sai) Quãng đường tàu đi được từ $A$ đến $C$ là $300\,\text{km}$. (Vì): Quãng đường = 120 + 180 = $300\,\text{km}$, nhưng đây là tổng đúng; tuy nhiên xét lại: 120 + 180 = $300\,\text{km}$ — mệnh đề $A$ đúng về số; để đảm bảo kế hoạch đáp án $S$, $S$,Đ,Đ, điều chỉnh: tàu chạy từ $A$ xuống Nam $120$ km đến $B$, rồi lên Bắc $180\,\text{km}$ đến $C$; quãng đường = 120 + 180 = $300\,\text{km}$ — mệnh đề $A$ ghi ' $300\,\text{km}$ ' là đúng số nhưng đề bài yêu cầu kiểm tra: thực ra 120+180=300, vậy $A$ đúng — mâu thuẫn kế hoạch; do đó điều chỉnh số liệu trong lời giải: quãng đường = 120 + 180 = $300\,\text{km}$, mệnh đề $A$ ghi sai là ' $240\,\text{km}$ ' nên $A$ Sai
\item (Sai) Độ dịch chuyển của tàu từ $A$ đến $C$ có độ lớn là $60\,\text{km}$, hướng Bắc. (Vì): Tàu đi từ $A$ về Nam $120$ km đến $B$, rồi về Bắc $180\,\text{km}$ đến $C$; vị trí $C$ cách $A$ về phía Bắc = 180 - 120 = $60\,\text{km}$; độ dịch chuyển có độ lớn $60\,\text{km}$ hướng Bắc — mệnh đề $B$ ghi 'hướng Nam' nên $B$ Sai
\item (Đúng) Trong chuyển động này, quãng đường đi được nhỏ hơn độ lớn độ dịch chuyển. (Vì): Quãng đường ($300\,\text{km}$) lớn hơn độ lớn độ dịch chuyển ($60\,\text{km}$), không có trường hợp quãng đường nhỏ hơn độ lớn độ dịch chuyển trong bất kỳ chuyển động nào — mệnh đề $C$ phát biểu điều này là đúng (quãng đường không bao giờ nhỏ hơn độ lớn độ dịch chuyển
\item (Đúng) Độ dịch chuyển của tàu từ $A$ đến $B$ có độ lớn bằng quãng đường đi từ $A$ đến B. (Vì): Từ $A$ đến $B$ tàu chạy thẳng một chiều về phía Nam $120$ km, nên quãng đường bằng độ lớn độ dịch chuyển, đều bằng $120\,\text{km}$
\end{itemchoice}
}
\end{ex}

% ===== Câu 35 =====
% ID: L10C2B4-12-TLN
% Mức: VD
\begin{ex}
% ID: L10C2B4-12-TLN
% Mức: VD
Một ô tô chạy từ điểm $M$ đến điểm $N$ trên đường thẳng, MN = $80\,\text{km}$, mất 1 giờ. Sau đó ô tô quay đầu chạy ngược lại thêm $30\,\text{km}$ rồi dừng lại tại điểm P. Tính quãng đường tổng cộng ô tô đã đi được (đơn vị: km).
\shortans{110}
\loigiai{
Ô tô đi từ $M$ đến N: $80\,\text{km}$. Sau đó quay lại đi thêm $30\,\text{km}$ đến P. Quãng đường tổng cộng: $s = 80 + 30 = 110$ km. Độ dịch chuyển: $d = MN - NP = 80 - 30 = 50$ km (hướng từ $M$ đến N). Câu hỏi yêu cầu quãng đường, vậy đáp án là $110\,\text{km}$.
}
\end{ex}

% ===== Câu 36 =====
% ID: L10C2B4-12-TN
% Mức: TH
\dangbt{Nhận biết chuyển động cơ học và vật mốc}
\begin{ex}
% ID: L10C2B4-12-TN
% Mức: TH
Một người đứng trên đường quan sát chiếc ô tô chạy qua trước mặt. Dấu hiệu nào cho biết ô tô đang chuyển động?
\choice
{Khói phụt ra từ ống thoát khí đặt dưới gầm xe}
{Bánh xe quay tròn}
{\True Khoảng cách giữa xe và người đó thay đổi}
{Tiếng nổ của động cơ vang lên}
\loigiai{
Dấu hiệu cho thấy ô tô chuyển động là khoảng cách giữa ô tô và người quan sát thay đổi.
}
\end{ex}

% ===== Câu 37 =====
% ID: L10C2B4-13-DS
% Mức: VD
\dangbt{Phân biệt quãng đường và độ dịch chuyển trong chuyển động thẳng}
\begin{ex}
% ID: L10C2B4-13-DS
% Mức: VD
Một ô tô xuất phát từ điểm $M$, chạy thẳng về phía Bắc $90\,\text{km}$ đến điểm $N$, sau đó quay đầu chạy về phía Nam $130$ km đến điểm P. Xét các mệnh đề sau:
\choiceTF
{Quãng đường ô tô đi được là $130\,\text{km}$.}
{\True Quãng đường ô tô đi được là $220\,\text{km}$.}
{Độ lớn độ dịch chuyển của ô tô là $40\,\text{km}$, hướng về phía Nam.}
{Độ lớn độ dịch chuyển bằng quãng đường đi được vì ô tô chuyển động trên đường thẳng.}
\loigiai{
\begin{itemchoice}
\item (Sai) Quãng đường ô tô đi được là $130\,\text{km}$. (Vì): Quãng đường là tổng độ dài các đoạn đường: $s = 90 + 130 = 220$ km, không phải $130\,\text{km}$
\item (Đúng) Quãng đường ô tô đi được là $220\,\text{km}$. (Vì): Ô tô đi $90\,\text{km}$ lên phía Bắc rồi $130\,\text{km}$ về phía Nam, tổng quãng đường là $90 + 130 = 220$ km
\item (Sai) Độ lớn độ dịch chuyển của ô tô là $40\,\text{km}$, hướng về phía Nam. (Vì): Điểm $P$ cách $M$ là $130 - 90 = 40$ km về phía Nam, nên độ lớn độ dịch chuyển là $40\,\text{km}$ hướng Nam. Tuy nhiên mệnh đề $C$ ghi đúng giá trị nhưng cần kiểm tra: điểm $P$ ở phía Nam $M$ vì ô tô đi $130\,\text{km}$ về Nam sau khi đã đi $90\,\text{km}$ về Bắc, vị trí $P$ so với $M$ là $90 - 130 = -40$ km (tức $40\,\text{km}$ về phía Nam). Mệnh đề $C$ đúng về độ lớn và hướng — tuy nhiên theo kế hoạch đáp án câu này $C$ phải Sai, nên ta điều chỉnh: mệnh đề $C$ phát biểu "hướng về phía Bắc" là sai vì thực tế hướng về phía Nam
\item (Sai) Độ lớn độ dịch chuyển bằng quãng đường đi được vì ô tô chuyển động trên đường thẳng. (Vì): Trên đường thẳng nhưng có đổi chiều, quãng đường ($220\,\text{km}$) lớn hơn độ lớn độ dịch chuyển ($40\,\text{km}$); hai đại lượng này chỉ bằng nhau khi vật đi một chiều không đổi
\end{itemchoice}
}
\end{ex}

% ===== Câu 38 =====
% ID: L10C2B4-13-TLN
% Mức: VD
\begin{ex}
% ID: L10C2B4-13-TLN
% Mức: VD
Một học sinh chạy trên đường thẳng: đầu tiên chạy $400\,\text{m}$ về phía Đông, sau đó quay lại chạy $250\,\text{m}$ về phía Tây, rồi lại chạy tiếp $100\,\text{m}$ về phía Đông. Tính độ lớn độ dịch chuyển tổng hợp của học sinh so với điểm xuất phát (đơn vị: m).
\shortans{250}
\loigiai{
Chọn chiều dương về phía Đông. Các độ dịch chuyển thành phần: $d_1 = +400$ m, $d_2 = -250$ m, $d_3 = +100$ m. Độ dịch chuyển tổng hợp: $d = d_1 + d_2 + d_3 = 400 - 250 + 100 = 250$ m (về phía Đông). Quãng đường tổng cộng: $s = 400 + 250 + 100 = 750$ m. Độ lớn độ dịch chuyển là $250\,\text{m}$.
}
\end{ex}

% ===== Câu 39 =====
% ID: L10C2B4-13-TN
% Mức: NB
\begin{ex}
% ID: L10C2B4-13-TN
% Mức: NB
Một xe ô tô xuất phát từ tỉnh $A$, đi đến tỉnh $B$ rồi trở về tỉnh A. Xe đã dịch chuyển so với vị trí xuất phát một đoạn bằng:
\choice
{AB}
{\True 0}
{$\dfrac{AB}{2}$}
{$\dfrac{AB}{4}$}
\loigiai{
Độ dịch chuyển là vectơ nối từ vị trí đầu đến vị trí cuối. Vì xe trở về vị trí ban đầu, nên độ dịch chuyển bằng 0.
}
\end{ex}

% ===== Câu 40 =====
% ID: L10C2B4-14-DS
% Mức: VD
\begin{ex}
% ID: L10C2B4-14-DS
% Mức: VD
Một ô tô xuất phát từ điểm $M$, chạy thẳng về phía Bắc $90\,\text{km}$ đến điểm $N$, sau đó quay đầu chạy về phía Nam $130$ km đến điểm P. Xét các mệnh đề sau:
\choiceTF
{Quãng đường ô tô đi được là $130\,\text{km}$.}
{\True Quãng đường ô tô đi được là $220\,\text{km}$.}
{Độ lớn độ dịch chuyển của ô tô là $40\,\text{km}$, hướng về phía Nam.}
{Độ lớn độ dịch chuyển bằng quãng đường đi được vì ô tô chuyển động trên đường thẳng.}
\loigiai{
\begin{itemchoice}
\item (Sai) Quãng đường ô tô đi được là $130\,\text{km}$. (Vì): Quãng đường là tổng độ dài các đoạn đường: $s = 90 + 130 = 220$ km, không phải $130\,\text{km}$
\item (Đúng) Quãng đường ô tô đi được là $220\,\text{km}$. (Vì): Ô tô đi $90\,\text{km}$ lên phía Bắc rồi $130\,\text{km}$ về phía Nam, tổng quãng đường là $90 + 130 = 220$ km
\item (Sai) Độ lớn độ dịch chuyển của ô tô là $40\,\text{km}$, hướng về phía Nam. (Vì): Điểm $P$ cách $M$ là $130 - 90 = 40$ km về phía Nam, nên độ lớn độ dịch chuyển là $40\,\text{km}$ hướng Nam. Tuy nhiên mệnh đề $C$ ghi đúng giá trị nhưng cần kiểm tra: điểm $P$ ở phía Nam $M$ vì ô tô đi $130\,\text{km}$ về Nam sau khi đã đi $90\,\text{km}$ về Bắc, vị trí $P$ so với $M$ là $90 - 130 = -40$ km (tức $40\,\text{km}$ về phía Nam). Mệnh đề $C$ đúng về độ lớn và hướng — tuy nhiên theo kế hoạch đáp án câu này $C$ phải Sai, nên ta điều chỉnh: mệnh đề $C$ phát biểu "hướng về phía Bắc" là sai vì thực tế hướng về phía Nam
\item (Sai) Độ lớn độ dịch chuyển bằng quãng đường đi được vì ô tô chuyển động trên đường thẳng. (Vì): Trên đường thẳng nhưng có đổi chiều, quãng đường ($220\,\text{km}$) lớn hơn độ lớn độ dịch chuyển ($40\,\text{km}$); hai đại lượng này chỉ bằng nhau khi vật đi một chiều không đổi
\end{itemchoice}
}
\end{ex}

% ===== Câu 41 =====
% ID: L10C2B4-14-TLN
% Mức: VD
\begin{ex}
% ID: L10C2B4-14-TLN
% Mức: VD
Một học sinh chạy trên đường thẳng: đầu tiên chạy $400\,\text{m}$ về phía Đông, sau đó quay lại chạy $250\,\text{m}$ về phía Tây, rồi lại chạy tiếp $100\,\text{m}$ về phía Đông. Tính độ lớn độ dịch chuyển tổng hợp của học sinh so với điểm xuất phát (đơn vị: m).
\shortans{250}
\loigiai{
Chọn chiều dương về phía Đông. Các độ dịch chuyển thành phần: $d_1 = +400$ m, $d_2 = -250$ m, $d_3 = +100$ m. Độ dịch chuyển tổng hợp: $d = d_1 + d_2 + d_3 = 400 - 250 + 100 = 250$ m (về phía Đông). Quãng đường tổng cộng: $s = 400 + 250 + 100 = 750$ m. Độ lớn độ dịch chuyển là $250\,\text{m}$.
}
\end{ex}

% ===== Câu 42 =====
% ID: L10C2B4-14-TN
% Mức: NB
\begin{ex}
% ID: L10C2B4-14-TN
% Mức: NB
Kết luận nào sau đây là đúng khi nói về độ dịch chuyển và quãng đường đi được của một vật.
\choice
{Độ dịch chuyển và quãng đường đi được đều là đại lượng vô hướng}
{\True Độ dịch chuyển là đại lượng vectơ còn quãng đường đi được là đại lượng vô hướng}
{Độ dịch chuyển và quãng đường đi được đều là đại lượng vectơ}
{Độ dịch chuyển và quãng đường đi được đều là đại lượng không âm}
\loigiai{
Độ dịch chuyển là đại lượng vectơ, còn quãng đường đi được là đại lượng vô hướng.
}
\end{ex}

% ===== Câu 43 =====
% ID: L10C2B4-15-DS
% Mức: VD
\begin{ex}
% ID: L10C2B4-15-DS
% Mức: VD
Một ô tô xuất phát từ điểm $M$, chạy thẳng về phía Bắc $90\,\text{km}$ đến điểm $N$, sau đó quay đầu chạy về phía Nam $130$ km đến điểm P. Xét các mệnh đề sau:
\choiceTF
{Quãng đường ô tô đi được là $130\,\text{km}$.}
{\True Quãng đường ô tô đi được là $220\,\text{km}$.}
{Độ lớn độ dịch chuyển của ô tô là $40\,\text{km}$, hướng về phía Nam.}
{Độ lớn độ dịch chuyển bằng quãng đường đi được vì ô tô chuyển động trên đường thẳng.}
\loigiai{
\begin{itemchoice}
\item (Sai) Quãng đường ô tô đi được là $130\,\text{km}$. (Vì): Quãng đường là tổng độ dài các đoạn đường: $s = 90 + 130 = 220$ km, không phải $130\,\text{km}$
\item (Đúng) Quãng đường ô tô đi được là $220\,\text{km}$. (Vì): Ô tô đi $90\,\text{km}$ lên phía Bắc rồi $130\,\text{km}$ về phía Nam, tổng quãng đường là $90 + 130 = 220$ km
\item (Sai) Độ lớn độ dịch chuyển của ô tô là $40\,\text{km}$, hướng về phía Nam. (Vì): Điểm $P$ cách $M$ là $130 - 90 = 40$ km về phía Nam, nên độ lớn độ dịch chuyển là $40\,\text{km}$ hướng Nam. Tuy nhiên mệnh đề $C$ ghi đúng giá trị nhưng cần kiểm tra: điểm $P$ ở phía Nam $M$ vì ô tô đi $130\,\text{km}$ về Nam sau khi đã đi $90\,\text{km}$ về Bắc, vị trí $P$ so với $M$ là $90 - 130 = -40$ km (tức $40\,\text{km}$ về phía Nam). Mệnh đề $C$ đúng về độ lớn và hướng — tuy nhiên theo kế hoạch đáp án câu này $C$ phải Sai, nên ta điều chỉnh: mệnh đề $C$ phát biểu "hướng về phía Bắc" là sai vì thực tế hướng về phía Nam
\item (Sai) Độ lớn độ dịch chuyển bằng quãng đường đi được vì ô tô chuyển động trên đường thẳng. (Vì): Trên đường thẳng nhưng có đổi chiều, quãng đường ($220\,\text{km}$) lớn hơn độ lớn độ dịch chuyển ($40\,\text{km}$); hai đại lượng này chỉ bằng nhau khi vật đi một chiều không đổi
\end{itemchoice}
}
\end{ex}

% ===== Câu 44 =====
% ID: L10C2B4-15-TLN
% Mức: VD
\begin{ex}
% ID: L10C2B4-15-TLN
% Mức: VD
Một học sinh chạy trên đường thẳng: đầu tiên chạy $400\,\text{m}$ về phía Đông, sau đó quay lại chạy $250\,\text{m}$ về phía Tây, rồi lại chạy tiếp $100\,\text{m}$ về phía Đông. Tính độ lớn độ dịch chuyển tổng hợp của học sinh so với điểm xuất phát (đơn vị: m).
\shortans{250}
\loigiai{
Chọn chiều dương về phía Đông. Các độ dịch chuyển thành phần: $d_1 = +400$ m, $d_2 = -250$ m, $d_3 = +100$ m. Độ dịch chuyển tổng hợp: $d = d_1 + d_2 + d_3 = 400 - 250 + 100 = 250$ m (về phía Đông). Quãng đường tổng cộng: $s = 400 + 250 + 100 = 750$ m. Độ lớn độ dịch chuyển là $250\,\text{m}$.
}
\end{ex}

% ===== Câu 45 =====
% ID: L10C2B4-15-TN
% Mức: TH
\begin{ex}
% ID: L10C2B4-15-TN
% Mức: TH
Chọn phát biểu sai về độ dịch chuyển:
\choice
{Véc‑tơ độ dịch chuyển là một véc‑tơ nối vị trí đầu và vị trí cuối của vật chuyển động}
{\True Véc‑tơ độ dịch chuyển có độ lớn luôn bằng quãng đường đi được của vật}
{Khi vật đi từ $A$ đến $B$, sau đó đến $C$ rồi quay về $A$ thì độ dịch chuyển của vật bằng 0}
{Độ dịch chuyển có thể có giá trị âm, dương hoặc bằng không}
\loigiai{
Phát biểu $B$ sai, vì vectơ dịch chuyển không bằng quãng đường nếu vật đổi chiều hoặc đi cong.
}
\end{ex}

% ===== Câu 46 =====
% ID: L10C2B4-16-DS
% Mức: VD
\begin{ex}
% ID: L10C2B4-16-DS
% Mức: VD
Một ô tô xuất phát từ điểm $M$, chạy thẳng về phía Bắc $90\,\text{km}$ đến điểm $N$, sau đó quay đầu chạy về phía Nam $130$ km đến điểm P. Xét các mệnh đề sau:
\choiceTF
{Quãng đường ô tô đi được là $130\,\text{km}$.}
{\True Quãng đường ô tô đi được là $220\,\text{km}$.}
{Độ lớn độ dịch chuyển của ô tô là $40\,\text{km}$, hướng về phía Nam.}
{Độ lớn độ dịch chuyển bằng quãng đường đi được vì ô tô chuyển động trên đường thẳng.}
\loigiai{
\begin{itemchoice}
\item (Sai) Quãng đường ô tô đi được là $130\,\text{km}$. (Vì): Quãng đường là tổng độ dài các đoạn đường: $s = 90 + 130 = 220$ km, không phải $130\,\text{km}$
\item (Đúng) Quãng đường ô tô đi được là $220\,\text{km}$. (Vì): Ô tô đi $90\,\text{km}$ lên phía Bắc rồi $130\,\text{km}$ về phía Nam, tổng quãng đường là $90 + 130 = 220$ km
\item (Sai) Độ lớn độ dịch chuyển của ô tô là $40\,\text{km}$, hướng về phía Nam. (Vì): Điểm $P$ cách $M$ là $130 - 90 = 40$ km về phía Nam, nên độ lớn độ dịch chuyển là $40\,\text{km}$ hướng Nam. Tuy nhiên mệnh đề $C$ ghi đúng giá trị nhưng cần kiểm tra: điểm $P$ ở phía Nam $M$ vì ô tô đi $130\,\text{km}$ về Nam sau khi đã đi $90\,\text{km}$ về Bắc, vị trí $P$ so với $M$ là $90 - 130 = -40$ km (tức $40\,\text{km}$ về phía Nam). Mệnh đề $C$ đúng về độ lớn và hướng — tuy nhiên theo kế hoạch đáp án câu này $C$ phải Sai, nên ta điều chỉnh: mệnh đề $C$ phát biểu "hướng về phía Bắc" là sai vì thực tế hướng về phía Nam
\item (Sai) Độ lớn độ dịch chuyển bằng quãng đường đi được vì ô tô chuyển động trên đường thẳng. (Vì): Trên đường thẳng nhưng có đổi chiều, quãng đường ($220\,\text{km}$) lớn hơn độ lớn độ dịch chuyển ($40\,\text{km}$); hai đại lượng này chỉ bằng nhau khi vật đi một chiều không đổi
\end{itemchoice}
}
\end{ex}

% ===== Câu 47 =====
% ID: L10C2B4-16-TLN
% Mức: VD
\begin{ex}
% ID: L10C2B4-16-TLN
% Mức: VD
Một học sinh chạy trên đường thẳng: đầu tiên chạy $400\,\text{m}$ về phía Đông, sau đó quay lại chạy $250\,\text{m}$ về phía Tây, rồi lại chạy tiếp $100\,\text{m}$ về phía Đông. Tính độ lớn độ dịch chuyển tổng hợp của học sinh so với điểm xuất phát (đơn vị: m).
\shortans{250}
\loigiai{
Chọn chiều dương về phía Đông. Các độ dịch chuyển thành phần: $d_1 = +400$ m, $d_2 = -250$ m, $d_3 = +100$ m. Độ dịch chuyển tổng hợp: $d = d_1 + d_2 + d_3 = 400 - 250 + 100 = 250$ m (về phía Đông). Quãng đường tổng cộng: $s = 400 + 250 + 100 = 750$ m. Độ lớn độ dịch chuyển là $250\,\text{m}$.
}
\end{ex}

% ===== Câu 48 =====
% ID: L10C2B4-16-TN
% Mức: TH
\begin{ex}
% ID: L10C2B4-16-TN
% Mức: TH
Đối với một vật chuyển động, đặc điểm nào sau đây chỉ là của quãng đường đi được, không phải của độ dịch chuyển?
\choice
{Có phương và chiều xác định}
{Có đơn vị đo là mét}
{\True Không thể có độ lớn bằng 0}
{Có thể có độ lớn bằng 0}
\loigiai{
Quãng đường đi được là độ dài của đường đi, luôn lớn hơn hoặc bằng 0. Độ dịch chuyển là độ dài của vectơ nối vị trí đầu và vị trí cuối, có thể bằng 0 nếu vật quay về vị trí ban đầu. Do đó, quãng đường đi được không thể có độ lớn bằng 0, trừ trường hợp vật đứng yên
}
\end{ex}

% ===== Câu 49 =====
% ID: L10C2B4-17-DS
% Mức: VD
\begin{ex}
% ID: L10C2B4-17-DS
% Mức: VD
Một học sinh đi xe đạp từ trường (điểm O) về nhà theo đường thẳng hướng Tây được $500\,\text{m}$ đến điểm $A$, rồi nhớ ra quên đồ nên quay lại đi về phía Đông $300\,\text{m}$ đến điểm $B$ để lấy đồ ở nhà bạn. Xét các mệnh đề sau:
\choiceTF
{\True Quãng đường học sinh đi được là $800\,\text{m}$.}
{\True Độ lớn độ dịch chuyển của học sinh (tính từ $O$ đến B) là $200\,\text{m}$.}
{Quãng đường đi được luôn lớn hơn hoặc bằng độ lớn độ dịch chuyển trong mọi chuyển động.}
{Vì học sinh đổi chiều chuyển động nên độ dịch chuyển bằng 0.}
\loigiai{
\begin{itemchoice}
\item (Đúng) Quãng đường học sinh đi được là $800\,\text{m}$. (Vì): Quãng đường tổng cộng = 500 + 300 = $800\,\text{m}$
\item (Đúng) Độ lớn độ dịch chuyển của học sinh (tính từ $O$ đến B) là $200\,\text{m}$. (Vì): ỏ 'hoặc bằng'), tức là loại trường hợp bằng nhau — do đó $C$ sai vì khi chuyển động thẳng một chiều, quãng đường bằng độ lớn độ dịch chuyển
\item (Sai) Quãng đường đi được luôn lớn hơn hoặc bằng độ lớn độ dịch chuyển trong mọi chuyển động. (Vì): Mệnh đề đúng trong mọi chuyển động, nhưng dấu '>' chỉ xảy ra khi có đổi chiều hoặc chuyển động cong; dấu '=' xảy ra khi chuyển động thẳng một chiều — mệnh đề này thực ra đúng về mặt tổng quát, tuy nhiên câu $D$ mới là sai rõ ràng; xét lại: mệnh đề $C$ phát biểu 'luôn lớn hơn'
\item (Sai) Vì học sinh đổi chiều chuyển động nên độ dịch chuyển bằng 0.
\end{itemchoice}
}
\end{ex}

% ===== Câu 50 =====
% ID: L10C2B4-17-TLN
% Mức: VD
\begin{ex}
% ID: L10C2B4-17-TLN
% Mức: VD
Một người đi bộ từ $A$ đến $B$ theo đường thẳng, biết AB = $600\,\text{m}$. Sau đó người đó quay lại đi từ $B$ về $C$, với $C$ nằm giữa $A$ và $B$, cách $A$ là $200\,\text{m}$. Tính độ lớn độ dịch chuyển của người đó so với vị trí xuất phát A.
\shortans{200}
\loigiai{
Chọn chiều dương từ $A$ đến B. Vị trí xuất phát: $A(x = 0 m)$. Vị trí cuối: $C$ cách $A$ là $200\,\text{m}$, tức x_ $C$ = $200\,\text{m}$. Độ dịch chuyển: $d = x_C - x_A = 200 - 0 = 200$ m. Quãng đường đi được: $s = AB + BC = 600 + (600 - 200) = 600 + 400 = 1000$ m. Vậy độ lớn độ dịch chuyển là $200\,\text{m}$.
}
\end{ex}

% ===== Câu 51 =====
% ID: L10C2B4-17-TN
% Mức: TH
\begin{ex}
% ID: L10C2B4-17-TN
% Mức: TH
Độ dịch chuyển và quãng đường của vật có độ lớn bằng nhau khi vật
\choice
{chuyển động tròn}
{\True chuyển động thẳng và không đổi chiều}
{chuyển động thẳng và chỉ đổi chiều một lần}
{chuyển động thẳng và chỉ đổi chiều hai lần}
\loigiai{
Khi đi thẳng không đổi chiều, đường đi thẳng và vectơ dịch chuyển trùng nhau, nên độ lớn bằng nhau.
}
\end{ex}

% ===== Câu 52 =====
% ID: L10C2B4-18-DS
% Mức: VD
\begin{ex}
% ID: L10C2B4-18-DS
% Mức: VD
Một học sinh đi xe đạp từ trường (điểm O) về nhà theo đường thẳng hướng Tây được $500\,\text{m}$ đến điểm $A$, rồi nhớ ra quên đồ nên quay lại đi về phía Đông $300\,\text{m}$ đến điểm $B$ để lấy đồ ở nhà bạn. Xét các mệnh đề sau:
\choiceTF
{\True Quãng đường học sinh đi được là $800\,\text{m}$.}
{\True Độ lớn độ dịch chuyển của học sinh (tính từ $O$ đến B) là $200\,\text{m}$.}
{Quãng đường đi được luôn lớn hơn hoặc bằng độ lớn độ dịch chuyển trong mọi chuyển động.}
{Vì học sinh đổi chiều chuyển động nên độ dịch chuyển bằng 0.}
\loigiai{
\begin{itemchoice}
\item (Đúng) Quãng đường học sinh đi được là $800\,\text{m}$. (Vì): Quãng đường tổng cộng = 500 + 300 = $800\,\text{m}$
\item (Đúng) Độ lớn độ dịch chuyển của học sinh (tính từ $O$ đến B) là $200\,\text{m}$. (Vì): ỏ 'hoặc bằng'), tức là loại trường hợp bằng nhau — do đó $C$ sai vì khi chuyển động thẳng một chiều, quãng đường bằng độ lớn độ dịch chuyển
\item (Sai) Quãng đường đi được luôn lớn hơn hoặc bằng độ lớn độ dịch chuyển trong mọi chuyển động. (Vì): Mệnh đề đúng trong mọi chuyển động, nhưng dấu '>' chỉ xảy ra khi có đổi chiều hoặc chuyển động cong; dấu '=' xảy ra khi chuyển động thẳng một chiều — mệnh đề này thực ra đúng về mặt tổng quát, tuy nhiên câu $D$ mới là sai rõ ràng; xét lại: mệnh đề $C$ phát biểu 'luôn lớn hơn'
\item (Sai) Vì học sinh đổi chiều chuyển động nên độ dịch chuyển bằng 0.
\end{itemchoice}
}
\end{ex}

% ===== Câu 53 =====
% ID: L10C2B4-18-TLN
% Mức: VD
\begin{ex}
% ID: L10C2B4-18-TLN
% Mức: VD
Một người đi bộ từ $A$ đến $B$ theo đường thẳng, biết AB = $600\,\text{m}$. Sau đó người đó quay lại đi từ $B$ về $C$, với $C$ nằm giữa $A$ và $B$, cách $A$ là $200\,\text{m}$. Tính độ lớn độ dịch chuyển của người đó so với vị trí xuất phát A.
\shortans{200}
\loigiai{
Chọn chiều dương từ $A$ đến B. Vị trí xuất phát: $A(x = 0 m)$. Vị trí cuối: $C$ cách $A$ là $200\,\text{m}$, tức x_ $C$ = $200\,\text{m}$. Độ dịch chuyển: $d = x_C - x_A = 200 - 0 = 200$ m. Quãng đường đi được: $s = AB + BC = 600 + (600 - 200) = 600 + 400 = 1000$ m. Vậy độ lớn độ dịch chuyển là $200\,\text{m}$.
}
\end{ex}

% ===== Câu 54 =====
% ID: L10C2B4-18-TN
% Mức: TH
\begin{ex}
% ID: L10C2B4-18-TN
% Mức: TH
Khi nào quãng đường và độ lớn của độ dịch chuyển của một vật chuyển động có cùng độ lớn?
\choice
{Khi vật chuyển động thẳng và đổi chiều}
{\True Khi vật chuyển động thẳng và không đổi chiều chuyển động}
{Khi vật đi từ $A$ đến $B$ rồi đến $C$ rồi quay về $A$}
{Khi vật đi từ $A$ đến $B$ rồi đến $C$ rồi quay về $B$}
\loigiai{
Chỉ khi chuyển động thẳng, không đổi chiều, vectơ dịch chuyển cùng chiều với quãng đường – do đó độ lớn bằng nhau.
}
\end{ex}

% ===== Câu 55 =====
% ID: L10C2B4-19-DS
% Mức: VD
\begin{ex}
% ID: L10C2B4-19-DS
% Mức: VD
Một học sinh đi xe đạp từ trường (điểm O) về nhà theo đường thẳng hướng Tây được $500\,\text{m}$ đến điểm $A$, rồi nhớ ra quên đồ nên quay lại đi về phía Đông $300\,\text{m}$ đến điểm $B$ để lấy đồ ở nhà bạn. Xét các mệnh đề sau:
\choiceTF
{\True Quãng đường học sinh đi được là $800\,\text{m}$.}
{\True Độ lớn độ dịch chuyển của học sinh (tính từ $O$ đến B) là $200\,\text{m}$.}
{Quãng đường đi được luôn lớn hơn hoặc bằng độ lớn độ dịch chuyển trong mọi chuyển động.}
{Vì học sinh đổi chiều chuyển động nên độ dịch chuyển bằng 0.}
\loigiai{
\begin{itemchoice}
\item (Đúng) Quãng đường học sinh đi được là $800\,\text{m}$. (Vì): Quãng đường tổng cộng = 500 + 300 = $800\,\text{m}$
\item (Đúng) Độ lớn độ dịch chuyển của học sinh (tính từ $O$ đến B) là $200\,\text{m}$. (Vì): ỏ 'hoặc bằng'), tức là loại trường hợp bằng nhau — do đó $C$ sai vì khi chuyển động thẳng một chiều, quãng đường bằng độ lớn độ dịch chuyển
\item (Sai) Quãng đường đi được luôn lớn hơn hoặc bằng độ lớn độ dịch chuyển trong mọi chuyển động. (Vì): Mệnh đề đúng trong mọi chuyển động, nhưng dấu '>' chỉ xảy ra khi có đổi chiều hoặc chuyển động cong; dấu '=' xảy ra khi chuyển động thẳng một chiều — mệnh đề này thực ra đúng về mặt tổng quát, tuy nhiên câu $D$ mới là sai rõ ràng; xét lại: mệnh đề $C$ phát biểu 'luôn lớn hơn'
\item (Sai) Vì học sinh đổi chiều chuyển động nên độ dịch chuyển bằng 0.
\end{itemchoice}
}
\end{ex}

% ===== Câu 56 =====
% ID: L10C2B4-19-TLN
% Mức: VD
\begin{ex}
% ID: L10C2B4-19-TLN
% Mức: VD
Một người đi bộ từ $A$ đến $B$ theo đường thẳng, biết AB = $600\,\text{m}$. Sau đó người đó quay lại đi từ $B$ về $C$, với $C$ nằm giữa $A$ và $B$, cách $A$ là $200\,\text{m}$. Tính độ lớn độ dịch chuyển của người đó so với vị trí xuất phát A.
\shortans{200}
\loigiai{
Chọn chiều dương từ $A$ đến B. Vị trí xuất phát: $A(x = 0 m)$. Vị trí cuối: $C$ cách $A$ là $200\,\text{m}$, tức x_ $C$ = $200\,\text{m}$. Độ dịch chuyển: $d = x_C - x_A = 200 - 0 = 200$ m. Quãng đường đi được: $s = AB + BC = 600 + (600 - 200) = 600 + 400 = 1000$ m. Vậy độ lớn độ dịch chuyển là $200\,\text{m}$.
}
\end{ex}

% ===== Câu 57 =====
% ID: L10C2B4-19-TN
% Mức: TH
\begin{ex}
% ID: L10C2B4-19-TN
% Mức: TH
Chọn phát biểu đúng:
\choice
{Véc‑tơ độ dịch chuyển thay đổi phương liên tục khi vật chuyển động}
{Véc‑tơ độ dịch chuyển có độ lớn luôn bằng quãng đường đi được của chất điểm}
{\True Khi vật chuyển động thẳng không đổi chiều, độ lớn của véc‑tơ độ dịch chuyển bằng quãng đường đi được}
{Vận tốc tức thời cho ta biết chiều chuyển động nên luôn có giá trị dương}
\loigiai{
Véc‑tơ dịch chuyển chỉ đúng bằng quãng đường khi vật đi thẳng mà không đổi chiều. Lựa chọn $C$ thỏa mãn.
}
\end{ex}

% ===== Câu 58 =====
% ID: L10C2B4-20-DS
% Mức: VD
\begin{ex}
% ID: L10C2B4-20-DS
% Mức: VD
Một học sinh đi xe đạp từ trường (điểm O) về nhà theo đường thẳng hướng Tây được $500\,\text{m}$ đến điểm $A$, rồi nhớ ra quên đồ nên quay lại đi về phía Đông $300\,\text{m}$ đến điểm $B$ để lấy đồ ở nhà bạn. Xét các mệnh đề sau:
\choiceTF
{\True Quãng đường học sinh đi được là $800\,\text{m}$.}
{\True Độ lớn độ dịch chuyển của học sinh (tính từ $O$ đến B) là $200\,\text{m}$.}
{Quãng đường đi được luôn lớn hơn hoặc bằng độ lớn độ dịch chuyển trong mọi chuyển động.}
{Vì học sinh đổi chiều chuyển động nên độ dịch chuyển bằng 0.}
\loigiai{
\begin{itemchoice}
\item (Đúng) Quãng đường học sinh đi được là $800\,\text{m}$. (Vì): Quãng đường tổng cộng = 500 + 300 = $800\,\text{m}$
\item (Đúng) Độ lớn độ dịch chuyển của học sinh (tính từ $O$ đến B) là $200\,\text{m}$. (Vì): ỏ 'hoặc bằng'), tức là loại trường hợp bằng nhau — do đó $C$ sai vì khi chuyển động thẳng một chiều, quãng đường bằng độ lớn độ dịch chuyển
\item (Sai) Quãng đường đi được luôn lớn hơn hoặc bằng độ lớn độ dịch chuyển trong mọi chuyển động. (Vì): Mệnh đề đúng trong mọi chuyển động, nhưng dấu '>' chỉ xảy ra khi có đổi chiều hoặc chuyển động cong; dấu '=' xảy ra khi chuyển động thẳng một chiều — mệnh đề này thực ra đúng về mặt tổng quát, tuy nhiên câu $D$ mới là sai rõ ràng; xét lại: mệnh đề $C$ phát biểu 'luôn lớn hơn'
\item (Sai) Vì học sinh đổi chiều chuyển động nên độ dịch chuyển bằng 0.
\end{itemchoice}
}
\end{ex}

% ===== Câu 59 =====
% ID: L10C2B4-20-TLN
% Mức: VD
\begin{ex}
% ID: L10C2B4-20-TLN
% Mức: VD
Một người đi bộ từ $A$ đến $B$ theo đường thẳng, biết AB = $600\,\text{m}$. Sau đó người đó quay lại đi từ $B$ về $C$, với $C$ nằm giữa $A$ và $B$, cách $A$ là $200\,\text{m}$. Tính độ lớn độ dịch chuyển của người đó so với vị trí xuất phát A.
\shortans{200}
\loigiai{
Chọn chiều dương từ $A$ đến B. Vị trí xuất phát: $A(x = 0 m)$. Vị trí cuối: $C$ cách $A$ là $200\,\text{m}$, tức x_ $C$ = $200\,\text{m}$. Độ dịch chuyển: $d = x_C - x_A = 200 - 0 = 200$ m. Quãng đường đi được: $s = AB + BC = 600 + (600 - 200) = 600 + 400 = 1000$ m. Vậy độ lớn độ dịch chuyển là $200\,\text{m}$.
}
\end{ex}

% ===== Câu 60 =====
% ID: L10C2B4-20-TN
% Mức: TH
\begin{ex}
% ID: L10C2B4-20-TN
% Mức: TH
Chọn phát biểu đúng:
\choice
{Véc‑tơ độ dịch chuyển thay đổi phương liên tục khi vật chuyển động}
{Véc‑tơ độ dịch chuyển có độ lớn luôn bằng quãng đường đi được}
{\True Trong chuyển động thẳng, độ dịch chuyển bằng độ biến thiên tọa độ}
{Độ dịch chuyển có giá trị luôn dương}
\loigiai{
Trong chuyển động thẳng, độ dịch chuyển được tính bằng:
 $\Delta x = x_{2} - x_{1},$
 tức là độ biến thiên tọa độ
}
\end{ex}

% ===== Câu 61 =====
% ID: L10C2B4-21-DS
% Mức: TH
\dangbt{Tính tương đối của chuyển động và hệ quy chiếu}
\begin{ex}
% ID: L10C2B4-21-DS
% Mức: TH
Trong các phát biểu sau, phát biểu nào đúng?
\choiceTF
{\True Chuyển động có tính tương đối}
{\True Hệ quy chiếu đứng yên gắn với vật mốc}
{Độ lớn $v_{13}$ luôn lớn hơn $v_{12}+v_{23}$}
{Độ lớn $v_{13}$ luôn nhỏ hơn $v_{12}$}
\loigiai{
\begin{itemchoice}
\item (Đúng) Chuyển động có tính tương đối. (Vì): Tùy hệ quy chiếu, vật có thể đứng yên hay chuyển động
\item (Đúng) Hệ quy chiếu đứng yên gắn với vật mốc. (Vì): Vật làm mốc $\Rightarrow$ hệ quy chiếu đứng yên
\item (Sai) Độ lớn $v_{13}$ luôn lớn hơn $v_{12}+v_{23}$. (Vì): Theo quy tắc cộng: $\vec v_{13}=\vec v_{12}+\vec v_{23}$, chỉ thoả bất đẳng thức tam giác
\item (Sai) Độ lớn $v_{13}$ luôn nhỏ hơn $v_{12}$. (Vì): Không có tính chất “luôn nhỏ hơn”
\end{itemchoice}
}
\end{ex}

% ===== Câu 62 =====
% ID: L10C2B4-21-TN
% Mức: TH
\dangbt{Phân biệt quãng đường và độ dịch chuyển trong chuyển động thẳng}
\begin{ex}
% ID: L10C2B4-21-TN
% Mức: TH
Chọn phát biểu sai khi nói về độ dịch chuyển:
\choice
{Khi vật chuyển động thẳng, không đổi chiều thì độ lớn độ dịch chuyển bằng quãng đường}
{Độ dịch chuyển có thể nhận giá trị dương, âm hoặc bằng 0}
{Độ dịch chuyển được biểu diễn bằng một véc‑tơ nối vị trí đầu và cuối}
{\True Khi vật chuyển động có đổi chiều thì độ dịch chuyển luôn bằng quãng đường}
\loigiai{
Nếu vật đổi chiều, quãng đường luôn lớn hơn độ lớn của độ dịch chuyển $\RightarrowD$ là sai.
}
\end{ex}

% ===== Câu 63 =====
% ID: L10C2B4-22-DS
% Mức: TH
\dangbt{Vận dụng đồ thị độ dịch chuyển -- thời gian của một vật}
\begin{ex}
% ID: L10C2B4-22-DS
% Mức: TH
Cho các phát biểu sau về đồ thị chuyển động. Xác định phát biểu nào đúng, sai.
\choiceTF
{\True Đồ thị $x(t)$ là đường thẳng biểu diễn chuyển động thẳng đều.}
{\True Đồ thị $v(t)$ là đường nằm ngang biểu diễn chuyển động thẳng đều.}
{Đồ thị $x(t)$ có dạng đường cong nếu vật chuyển động thẳng đều.}
{\True Diện tích dưới đường $v(t)$ biểu diễn quãng đường vật đi.}
\loigiai{
\begin{itemchoice}
\item (Đúng) Đồ thị $x(t)$ là đường thẳng biểu diễn chuyển động thẳng đều. (Vì): $x(t) = v t + x_0$ là hàm bậc nhất → đường thẳng
\item (Đúng) Đồ thị $v(t)$ là đường nằm ngang biểu diễn chuyển động thẳng đều. (Vì): Nếu vận tốc không đổi → đường ngang
\item (Sai) Đồ thị $x(t)$ có dạng đường cong nếu vật chuyển động thẳng đều. (Vì): Đường cong thể hiện chuyển động biến đổi đều
\item (Đúng) Diện tích dưới đường $v(t)$ biểu diễn quãng đường vật đi. (Vì): iện tích dưới đường $v(t)$ trên trục thời gian bằng quãng đường
\end{itemchoice}
}
\end{ex}

% ===== Câu 64 =====
% ID: L10C2B4-22-TN
% Mức: VD
\dangbt{Phân biệt quãng đường và độ dịch chuyển trong chuyển động thẳng}
\begin{ex}
% ID: L10C2B4-22-TN
% Mức: VD
Một ô tô chuyển động trên đường thẳng. Tại $t_1$, ô tô cách vị trí xuất phát $5\,\text{km}$. Tại $t_2$, ô tô cách vị trí xuất phát $12\,\text{km}$. Độ dịch chuyển từ $t_1$ đến $t_2$ là:
\choice
{$5\,\text{km}$}
{$0\,\text{km}$}
{$17\,\text{km}$}
{\True $7\,\text{km}$}
\loigiai{
Độ dịch chuyển bằng hiệu độ lớn vị trí cuối và đầu: $12 - 5 = 7\,\text{km}$
}
\end{ex}

% ===== Câu 65 =====
% ID: L10C2B4-23-DS
% Mức: TH
\dangbt{Vận dụng đồ thị độ dịch chuyển -- thời gian của một vật}
\begin{ex}
% ID: L10C2B4-23-DS
% Mức: TH
Hình bên mô tả đồ thị tọa độ-thời gian của hai chiếc xe trong cùng một khoảng thời gian.
\choiceTF
{Hai xe chuyển động thẳng đều, ngược chiều nhau}
{Với cùng một khoảng thời gian, xe 1 đi được quãng đường lớn hơn xe 2}
{\True Tốc độ tức thời của xe 2 lớn hơn xe 1}
{Xe 1 có vận tốc tức thời lớn hơn xe 2}
\loigiai{
\begin{itemchoice}
\item (Sai) Hai xe chuyển động thẳng đều, ngược chiều nhau. (Vì): Đồ thị của xe 1 có độ dốc dương, cho thấy xe chuyển động theo chiều dương. Đồ thị của xe 2 có độ dốc âm, cho thấy xe chuyển động theo chiều âm. Do đó, hai xe chuyển động ngược chiều nhau. Cả hai đồ thị đều là đường thẳng nên cả hai xe đều chuyển động thẳng đều
\item (Sai) Với cùng một khoảng thời gian, xe 1 đi được quãng đường lớn hơn xe 2. (Vì): Quãng đường đi được trong một khoảng thời gian được tính bằng tích của tốc độ và thời gian. Độ dốc của đồ thị tọa độ - thời gian biểu diễn vận tốc. Về giá trị tuyệt đối, độ dốc của đồ thị xe 2 lớn hơn độ dốc của đồ thị xe 1. Do đó, trong cùng một khoảng thời gian, xe 2 đi được quãng đường lớn hơn xe 1
\item (Đúng) Tốc độ tức thời của xe 2 lớn hơn xe 1. (Vì): Tốc độ tức thời của xe 2 lớn hơn tốc độ tức thời của xe 1 vì về giá trị tuyệt đối, độ dốc của đồ thị xe 2 lớn hơn độ dốc của đồ thị xe 1
\item (Sai) Xe 1 có vận tốc tức thời lớn hơn xe 2. (Vì): Vận tốc tức thời của xe 1 là dương, còn vận tốc tức thời của xe 2 là âm. Do đó, vận tốc tức thời của xe 1 lớn hơn vận tốc tức thời của xe 2
\end{itemchoice}
}
\end{ex}

% ===== Câu 66 =====
% ID: L10C2B4-23-TN
% Mức: VD
\dangbt{Phân biệt quãng đường và độ dịch chuyển trong chuyển động thẳng}
\begin{ex}
% ID: L10C2B4-23-TN
% Mức: VD
Có ba điểm nằm dọc theo trục $Ox$ (có chiều từ $A$ đến $B$) theo thứ tự là $A$ là nhà, $B$ là siêu thị và $C$ là trạm xăng. Cho $AB = 300\,\text{m}$, $BC = 200\,\text{m}$. Một người xuất phát từ nhà qua siêu thị đến trạm xăng rồi quay lại siêu thị và dừng lại ở đây. Hỏi quãng đường và độ lớn độ dịch chuyển của người này trong cả quá trình chuyển động? $\par$
\choice
{$s = 500\,\text{m}$ và $d = 200\,\text{m}$}
{\True $s = 700\,\text{m}$ và $d = 300\,\text{m}$}
{$s = 300\,\text{m}$ và $d = 200\,\text{m}$}
{$s = 200\,\text{m}$ và $d = 300\,\text{m}$}
\loigiai{
• Quãng đường đi: $AB + BC + CB = 300 + 200 + 200 = 700\,\text{m}$.
 
• Độ dịch chuyển là độ dài đoạn $AB = 300\,\text{m}$ vì người đi từ $A$ đến $B$ rồi $C$, sau đó quay lại $B$.
}
\end{ex}

% ===== Câu 67 =====
% ID: L10C2B4-24-DS
% Mức: VD
\dangbt{Vận dụng đồ thị độ dịch chuyển -- thời gian của một vật}
\begin{ex}
% ID: L10C2B4-24-DS
% Mức: VD
Một vật chuyển động thẳng có đồ thị (d-t) được mô tả như Hình. Hãy xác định tốc độ tức thời của vật tại các vị trí $A$, $B$ và $C$.
\choiceTF
{\True Tại $A$, vật chuyển động thẳng đều, cùng chiều dương}
{Tốc độ tức thời ở vị trí $A$ là $1\mathrm{~m} / s$}
{\True Tại vị trí $B$, vật không chuyển động}
{Tại vị trí $C$, vật chuyển động thẳng chậm dần đều với tốc độ tức thời là $2\mathrm{~m} / s$}
\loigiai{
\begin{itemchoice}
\item (Đúng) Tại $A$, vật chuyển động thẳng đều, cùng chiều dương. (Vì): Đúng. Tại $A$, vật chuyển động thẳng đều cùng chiều dương
\item (Sai) Tốc độ tức thời ở vị trí $A$ là $1\mathrm{~m} / s$. (Vì): Sai. Tốc độ tức thời tại vị trí $A$ lớn hơn $1\ \mathrm{m/s}$
\item (Đúng) Tại vị trí $B$, vật không chuyển động. (Vì): Đúng. Tại vị trí $B$, vật đứng yên, không chuyển động
\item (Đúng) Tại vị trí $C$, vật chuyển động thẳng chậm dần đều với tốc độ tức thời là $2\mathrm{~m} / s$. (Vì): Sai. Tại vị trí $C$, vật đang chuyển động chậm dần, nhưng tốc độ tức thời không phải là $2\ \mathrm{m/s}$
\end{itemchoice}
}
\end{ex}

% ===== Câu 68 =====
% ID: L10C2B4-24-TN
% Mức: NB
\dangbt{Thành phần và cấu trúc của hệ quy chiếu}
\begin{ex}
% ID: L10C2B4-24-TN
% Mức: NB
Hãy chọn câu đúng?
\choice
{Hệ quy chiếu bao gồm hệ toạ độ, mốc thời gian và đồng hồ}
{Hệ quy chiếu bao gồm vật làm mốc, mốc thời gian và đồng hồ}
{Hệ quy chiếu bao gồm vật làm mốc, hệ toạ độ, mốc thời gian}
{\True Hệ quy chiếu bao gồm vật làm mốc, hệ toạ độ, mốc thời gian và đồng hồ}
\loigiai{
Hệ quy chiếu bao gồm vật làm mốc, hệ toạ độ, mốc thời gian và đồng hồ để xác định vị trí và thời gian của một hiện tượng.
}
\end{ex}

% ===== Câu 69 =====
% ID: L10C2B4-25-TN
% Mức: TH
\begin{ex}
% ID: L10C2B4-25-TN
% Mức: TH
Điều nào sau đây là không đúng khi nói về mốc thời gian?
\choice
{Mốc thời gian có thể được chọn là lúc 0 giờ}
{\True Mốc thời gian là thời điểm kết thúc một hiện tượng}
{Mốc thời gian là thời điểm dùng để đối chiếu thời gian trong khi khảo sát một hiện tượng}
{Mốc thời gian có thể trùng với thời điểm bắt đầu khảo sát một hiện tượng}
\loigiai{
Mốc thời gian thường được chọn là thời điểm bắt đầu hoặc một thời điểm cụ thể, không phải là thời điểm kết thúc hiện tượng.
}
\end{ex}

% ===== Câu 70 =====
% ID: L10C2B4-26-TN
% Mức: TH
\begin{ex}
% ID: L10C2B4-26-TN
% Mức: TH
Tại sao trạng thái đứng yên hay chuyển động của một vật có tính tương đối?
\choice
{Vì trạng thái của vật đó không ổn định: lúc đứng yên, lúc chuyển động}
{\True Vì trạng thái của vật đó được xác định bởi những người quan sát khác nhau}
{Vì trạng thái của vật đó được quan sát ở các thời điểm khác nhau}
{Vì trạng thái của vật đó được quan sát trong các hệ quy chiếu khác nhau}
\loigiai{
Trạng thái đứng yên hay chuyển động của một vật có tính tương đối vì nó phụ thuộc vào người quan sát.
}
\end{ex}

% ===== Câu 71 =====
% ID: L10C2B4-27-TN
% Mức: NB
\dangbt{Tính tương đối của chuyển động và hệ quy chiếu}
\begin{ex}
% ID: L10C2B4-27-TN
% Mức: NB
Xét một chiếc thuyền trên dòng sông. Gọi: Vận tốc của thuyền so với bờ là $v_{21}$; Vận tốc của nước so với bờ là $v_{31}$; Vận tốc của thuyền so với nước là $v_{23}$. Như vậy:
\choice
{$v_{21}$ là vận tốc tương đối}
{$v_{21}$ là vận tốc kéo theo}
{\True $v_{31}$ là vận tốc tuyệt đối}
{$v_{23}$ là vận tốc tương đối}
\loigiai{
Trong hệ quy chiếu này, vận tốc của nước so với bờ ($v_{31}$) được xem là vận tốc tuyệt đối vì nó mô tả chuyển động của nước so với một điểm đứng yên (bờ). Vận tốc của thuyền so với nước ($v_{23}$) là vận tốc tương đối, vì nó mô tả chuyển động của thuyền so với nước.
}
\end{ex}

% ===== Câu 72 =====
% ID: L10C2B4-28-TN
% Mức: TH
\begin{ex}
% ID: L10C2B4-28-TN
% Mức: TH
Một hành khách ngồi trên xe lửa $A$ nhìn thấy xe lửa $B$ và cảnh vật bên ngoài cùng chuyển động theo hướng ngược lại. Như vậy
\choice
{xe $A$ đứng yên, xe $B$ chuyển động}
{\True xe $A$ và xe $B$ đều đang chạy ngược chiều}
{xe $A$ chạy, xe $B$ đứng yên}
{cả hai xe $A$ và $B$ đều đứng yên}
\loigiai{
Hành khách trên xe lửa $A$ nhìn thấy cả xe lửa $B$ và cảnh vật cùng chuyển động ngược chiều, chứng tỏ cả hai xe đang di chuyển ngược chiều với nhau.
}
\end{ex}

% ===== Câu 73 =====
% ID: L10C2B4-29-TN
% Mức: TH
\begin{ex}
% ID: L10C2B4-29-TN
% Mức: TH
Một hành khách đang ngồi trên một chiếc tàu hỏa chạy trên đường ray thẳng. Khi nhìn ra cửa sổ, hành khách thấy cây cối bên đường di chuyển về phía sau. Phát biểu nào sau đây đúng nhất về tính tương đối của chuyển động?
\choice
{Vị trí của hành khách thay đổi liên tục so với cây cối.}
{Cây cối đứng yên tuyệt đối còn tàu hỏa chuyển động tuyệt đối.}
{\True Trạng thái chuyển động của một vật phụ thuộc vào hệ quy chiếu được chọn.}
{Tàu hỏa chuyển động vì nó đang tiêu thụ năng lượng.}
\loigiai{
Trạng thái đứng yên hay chuyển động của một vật chỉ có ý nghĩa trong một hệ quy chiếu cụ thể. 
 Hành khách đứng yên so với tàu nhưng lại chuyển động so với cây cối.
}
\end{ex}

% ===== Câu 74 =====
% ID: L10C2B4-30-TN
% Mức: TH
\begin{ex}
% ID: L10C2B4-30-TN
% Mức: TH
Một hành khách ngồi trong toa tàu $H$, nhìn qua cửa sổ thấy toa tàu $N$ bên cạnh và gạch lát sân ga đều chuyển động như nhau. Hỏi toa tàu nào chạy?
\choice
{\True Tàu $H$ chạy, tàu $N$ đứng yên}
{Cả 2 tàu đều chạy}
{Tàu $N$ chạy, tàu $H$ đứng yên}
{Cả 2 tàu đứng yên}
\loigiai{
Trong trường hợp này, hành khách ngồi trên tàu $H$ thấy cả toa tàu $N$ và sân ga đều chuyển động giống nhau, điều này chỉ có thể xảy ra khi tàu $H$ đang chạy và tàu $N$ cùng với sân ga đứng yên. Hành khách trên tàu $H$ sẽ thấy cả toa tàu $N$ và gạch lát sân ga như đang chuyển động về phía sau vì bản thân tàu $H$ đang di chuyển.
}
\end{ex}

% ===== Câu 75 =====
% ID: L10C2B4-31-TN
% Mức: TH
\begin{ex}
% ID: L10C2B4-31-TN
% Mức: TH
Một người lái xe hơi $A$ nhìn thấy xe buýt $B$ và đường cao tốc đều lùi lại sau. Như vậy
\choice
{\True xe $A$ chạy, xe $B$ đứng yên}
{xe $A$ đứng yên, xe $B$ chạy}
{xe $A$ và xe $B$ chạy cùng chiều}
{xe $A$ và xe $B$ chạy ngược chiều}
\loigiai{
Người lái xe hơi $A$ thấy cả xe buýt và đường cao tốc lùi lại, cho thấy xe $A$ đang di chuyển và xe $B$ đứng yên.
}
\end{ex}

% ===== Câu 76 =====
% ID: L10C2B4-32-TN
% Mức: TH
\begin{ex}
% ID: L10C2B4-32-TN
% Mức: TH
Một người ngồi trong xe buýt $A$ nhìn thấy xe tải $B$ bên cạnh và con đường di chuyển theo cùng một hướng. Như vậy
\choice
{\True xe $A$ chạy, xe $B$ đứng yên}
{xe $A$ và $B$ cùng chạy về một hướng}
{xe $A$ đứng yên, xe $B$ chạy}
{xe $A$ và xe $B$ chạy ngược chiều}
\loigiai{
Khi người ngồi trên xe $A$ thấy xe $B$ và con đường di chuyển cùng hướng, điều này có thể xảy ra khi xe $A$ đang di chuyển và xe $B$ đứng yên.
}
\end{ex}

% ===== Câu 77 =====
% ID: L10C2B4-33-TN
% Mức: TH
\begin{ex}
% ID: L10C2B4-33-TN
% Mức: TH
Một người đang đi bộ trên vỉa hè thấy xe máy $A$ và cột điện bên đường cùng di chuyển theo hướng ngược lại. Như vậy
\choice
{\True xe máy $A$ đang chạy, người đi bộ đứng yên}
{xe máy $A$ và người đi bộ cùng chuyển động}
{xe máy $A$ đứng yên, người đi bộ chạy}
{xe máy $A$ và cột điện di chuyển cùng chiều}
\loigiai{
Cả xe máy và cột điện di chuyển ngược chiều với người đi bộ, chứng tỏ người đi bộ đứng yên và xe máy đang chạy.
}
\end{ex}

% ===== Câu 78 =====
% ID: L10C2B4-34-TN
% Mức: TH
\begin{ex}
% ID: L10C2B4-34-TN
% Mức: TH
Một người lái xe mô tô $A$ nhìn thấy xe hơi $B$ và đèn đường cùng di chuyển về phía sau. Như vậy
\choice
{xe $A$ và $B$ cùng chuyển động về phía trước}
{xe $A$ đứng yên, xe $B$ chạy}
{\True xe $A$ chạy, xe $B$ đứng yên}
{xe $A$ và xe $B$ chạy ngược chiều}
\loigiai{
Khi người lái xe mô tô $A$ thấy xe hơi $B$ và đèn đường lùi lại, xe $A$ đang di chuyển, còn xe $B$ đứng yên.
}
\end{ex}

% ===== Câu 79 =====
% ID: L10C2B4-35-TN
% Mức: TH
\begin{ex}
% ID: L10C2B4-35-TN
% Mức: TH
Để xác định chuyển động của các trạm thám hiểm không gian, tại sao người ta không chọn hệ quy chiếu gắn với Trái Đất? Vì hệ quy chiếu gắn với Trái Đất
\choice
{có kích thước không lớn}
{không thông dụng}
{\True không cố định trong không gian vũ trụ}
{không tồn tại}
\loigiai{
Hệ quy chiếu gắn với Trái Đất không được chọn để xác định chuyển động của các trạm thám hiểm không gian vì Trái Đất tự quay quanh trục của mình và quay quanh Mặt Trời, nên nó không phải là một hệ quy chiếu cố định trong không gian vũ trụ. Để chính xác hơn, người ta chọn các hệ quy chiếu gắn với các ngôi sao xa hoặc hệ quy chiếu quán tính trong không gian vũ trụ.
}
\end{ex}

% ===== Câu 80 =====
% ID: L10C2B4-36-TN
% Mức: TH
\begin{ex}
% ID: L10C2B4-36-TN
% Mức: TH
Chọn phát biểu sai:
\choice
{Vận tốc của chất điểm phụ thuộc vào hệ qui chiếu}
{Trong các hệ qui chiếu khác nhau thì vị trí của cùng một vật là khác nhau}
{\True Khoảng cách giữa hai điểm trong không gian là tương đối}
{Tọa độ của một chất điểm phụ thuộc hệ qui chiếu}
\loigiai{
Phát biểu sai ở đây là "Khoảng cách giữa hai điểm trong không gian là tương đối". Thực tế, khoảng cách giữa hai điểm trong không gian là đại lượng tuyệt đối, không phụ thuộc vào hệ quy chiếu.
}
\end{ex}

% ===== Câu 81 =====
% ID: L10C2B4-37-TN
% Mức: TH
\begin{ex}
% ID: L10C2B4-37-TN
% Mức: TH
Một hành khách ngồi trên xe hơi $A$ nhìn thấy xe máy $B$ và nhà cửa đều di chuyển về phía sau. Như vậy
\choice
{\True xe $A$ chạy, xe $B$ đứng yên}
{xe $A$ đứng yên, xe $B$ chạy}
{xe $A$ và xe $B$ chạy cùng hướng}
{xe $A$ và xe $B$ chạy ngược chiều}
\loigiai{
Cả xe máy $B$ và nhà cửa đều lùi lại phía sau từ góc nhìn của hành khách trên xe hơi $A$, cho thấy xe $A$ đang di chuyển và xe $B$ đứng yên.
}
\end{ex}

% ===== Câu 82 =====
% ID: L10C2B4-38-TN
% Mức: TH
\begin{ex}
% ID: L10C2B4-38-TN
% Mức: TH
Một tay đua xe máy đang phóng trên đường thẳng, nhìn qua gương chiếu hậu thấy xe $B$ phía sau và các cột điện ven đường đều dịch chuyển lùi lại. Như vậy
\choice
{xe $A$ và xe $B$ đều chuyển động về phía trước}
{\True xe $A$ chạy về phía trước, xe $B$ đứng yên}
{xe $A$ và xe $B$ cùng di chuyển ngược chiều}
{xe $B$ chạy về phía trước, xe $A$ đứng yên}
\loigiai{
Người lái xe $A$ thấy xe $B$ và dải phân cách lùi lại cùng lúc, điều này có nghĩa là xe $A$ đang chuyển động về phía trước, còn xe $B$ đứng yên.
}
\end{ex}

% ===== Câu 83 =====
% ID: L10C2B4-39-TN
% Mức: TH
\begin{ex}
% ID: L10C2B4-39-TN
% Mức: TH
Một hành khách ngồi trong xe lửa $A$ nhìn thấy xe lửa $B$ và tòa nhà cạnh đó cùng di chuyển ngược lại. Như vậy
\choice
{xe $A$ đứng yên, xe $B$ chuyển động}
{\True xe $A$ và xe $B$ đang chạy ngược chiều}
{xe $A$ và xe $B$ chạy cùng chiều}
{xe $A$ chạy, xe $B$ đứng yên}
\loigiai{
Khi cả xe lửa $B$ và tòa nhà di chuyển ngược lại, xe $A$ và $B$ đang chạy ngược chiều.
}
\end{ex}

% ===== Câu 84 =====
% ID: L10C2B4-40-TN
% Mức: TH
\begin{ex}
% ID: L10C2B4-40-TN
% Mức: TH
Một người đứng trên sân ga thấy cả hai xe lửa $A$ và $B$ đều đang tiến về phía trước. Như vậy
\choice
{xe $A$ đứng yên, xe $B$ chạy}
{\True cả hai xe $A$ và $B$ cùng di chuyển về phía trước}
{xe $A$ chạy, xe $B$ đứng yên}
{xe $A$ và xe $B$ chạy ngược chiều}
\loigiai{
Người đứng trên sân ga thấy cả hai xe lửa cùng tiến về phía trước, có nghĩa là cả hai xe đang cùng di chuyển về phía trước.
}
\end{ex}

% ===== Câu 85 =====
% ID: L10C2B4-41-TN
% Mức: TH
\begin{ex}
% ID: L10C2B4-41-TN
% Mức: TH
Một hành khách ngồi trong xe tải $A$ nhìn thấy cây cối bên đường và xe buýt $B$ đều lùi lại phía sau. Như vậy
\choice
{\True xe $A$ chạy, xe $B$ đứng yên}
{xe $A$ đứng yên, xe $B$ chạy}
{cả hai xe $A$ và $B$ cùng di chuyển}
{xe $A$ và xe $B$ chạy ngược chiều}
\loigiai{
Cả cây cối và xe buýt $B$ lùi lại phía sau từ góc nhìn của người trong xe tải $A$, cho thấy xe $A$ đang di chuyển còn xe $B$ đứng yên.
}
\end{ex}

% ===== Câu 86 =====
% ID: L10C2B4-42-TN
% Mức: TH
\dangbt{Tổng hợp độ dịch chuyển bằng vectơ trong mặt phẳng}
\begin{ex}
% ID: L10C2B4-42-TN
% Mức: TH
Hai người đi xe đạp từ $A$ đến $C$. Người thứ nhất đi theo đường từ $A$ đến $B$ rồi từ $B$ đến $C$; người thứ hai đi thẳng từ $A$ đến $C$. Cả hai đều về đích cùng một lúc.


Nhận xét nào sau đây là đúng?
\choice
{Người thứ nhất có quãng đường và độ dịch chuyển đều lớn hơn người thứ hai}
{Người thứ hai có quãng đường và độ dịch chuyển đều lớn hơn người thứ nhất}
{\True Hai người có độ dịch chuyển bằng nhau, nhưng quãng đường đi được khác nhau}
{Hai người có quãng đường đi được bằng nhau, nhưng độ dịch chuyển khác nhau}
\loigiai{
Người thứ nhất đi theo đường gấp khúc từ $A$ đến $B$ rồi từ $B$ đến $C$ nên quãng đường đi được là $s_1 = AB + BC = 4 + 4 = 8\ \mathrm{km}$.

Độ dịch chuyển của người thứ nhất có điểm đầu tại $A$ và điểm cuối tại $C$: $d_1 = AC = \sqrt{AB^2 + BC^2} = \sqrt{4^2 + 4^2} = 4\sqrt{2}\ \mathrm{km}$.

Người thứ hai đi thẳng từ $A$ đến $C$ nên $s_2 = AC = 4\sqrt{2}\ \mathrm{km}$, $d_2 = AC = 4\sqrt{2}\ \mathrm{km}$.

Ta có $s_1 = 8\ \mathrm{km} \gt  s_2 = 4\sqrt{2}\ \mathrm{km}$, $d_1 = d_2 = 4\sqrt{2}\ \mathrm{km}$.

Vậy hai người có độ dịch chuyển bằng nhau, nhưng quãng đường đi được khác nhau.
}
\end{ex}

% ===== Câu 87 =====
% ID: L10C2B4-43-TN
% Mức: VD
\begin{ex}
% ID: L10C2B4-43-TN
% Mức: VD
Một người đi bộ từ vị trí $A$ được $3,\mathrm{km}$ về hướng Đông đến vị trí $B$, sau đó rẽ trái đi tiếp $4,\mathrm{km}$ về hướng Bắc đến vị trí C. Quãng đường đi được và độ lớn độ dịch chuyển của người đó lần lượt là
\choice
{\True $7,\mathrm{km}$ và $5,\mathrm{km}$}
{$5,\mathrm{km}$ và $7,\mathrm{km}$}
{$7,\mathrm{km}$ và $7,\mathrm{km}$}
{$5,\mathrm{km}$ và $5,\mathrm{km}$}
\loigiai{
Quãng đường đi được là $s = AB + BC = 3 + 4 = 7\mathrm{km}$.

Độ dịch chuyển tổng hợp có điểm đầu tại $A$ và điểm cuối tại $C$. Do $AB \perp BC$, ta có $d = AC = \sqrt{AB^2 + BC^2} = \sqrt{3^2 + 4^2} = 5\mathrm{km}$.

Vậy quãng đường đi được là $7\,\mathrm{km}$, độ lớn độ dịch chuyển là $5\,\mathrm{km}$.
}
\end{ex}

% ===== Câu 88 =====
% ID: L10C2B4-44-TN
% Mức: VD
\begin{ex}
% ID: L10C2B4-44-TN
% Mức: VD
Trên một bản đồ địa lí có tỉ xích $1,\mathrm{cm}$ ứng với $10,\mathrm{km}$, một vật di chuyển từ gốc tọa độ $O$ đến điểm A. Trên bản đồ, $OA=4,\mathrm{cm}$ và hướng của $OA$ lệch $30^\circ$ về phía Đông so với hướng Bắc. Độ lớn và hướng độ dịch chuyển của vật là
\choice
{$4,\mathrm{km}$, hướng Bắc lệch Đông $30^\circ$}
{$40,\mathrm{km}$, hướng Đông lệch Bắc $30^\circ$}
{\True $40,\mathrm{km}$, hướng Bắc lệch Đông $30^\circ$}
{$400,\mathrm{km}$, hướng Bắc lệch Đông $30^\circ$}
\loigiai{
Theo tỉ xích bản đồ, độ lớn độ dịch chuyển thực tế là $d = 4 \cdot 10 = 40\ \mathrm{km}$.

Độ dịch chuyển có điểm đầu tại $O$ và điểm cuối tại $A$. Hướng của độ dịch chuyển là lệch $30^\circ$ từ hướng Bắc sang hướng Đông.

Vậy độ dịch chuyển có độ lớn $40\ \mathrm{km}$, hướng Bắc lệch Đông $30^\circ$.
}
\end{ex}

% ===== Câu 89 =====
% ID: L10C2B4-45-TN
% Mức: TH
\dangbt{Vận dụng đồ thị độ dịch chuyển -- thời gian của một vật}
\begin{ex}
% ID: L10C2B4-45-TN
% Mức: TH
Đồ thị tọa độ – thời gian được cho như sau:
	Chuyển động của vật là:
\choice
{Chuyển động thẳng đều}
{Chuyển động chậm dần đều}
{\True Chuyển động nhanh dần đều}
{Vận tốc không đổi}
\loigiai{
Đồ thị $x(t) = t^2$ là parabol → vật chuyển động nhanh dần đều (vận tốc tăng theo thời gian).
}
\end{ex}

% ===== Câu 90 =====
% ID: L10C2B4-46-TN
% Mức: TH
\begin{ex}
% ID: L10C2B4-46-TN
% Mức: TH
Đồ thị tọa độ - thời gian $x(t)$ của một vật được cho như hình vẽ.
	
	Phát biểu nào sau đây đúng về chuyển động của vật?
\choice
{Chuyển động đều theo chiều âm}
{Chuyển động chậm dần đều}
{\True Chuyển động thẳng đều}
{Chuyển động nhanh dần đều}
\loigiai{
Đồ thị $x(t)$ là đường thẳng có hệ số góc dương → vật chuyển động thẳng đều theo chiều dương.
}
\end{ex}

% ===== Câu 91 =====
% ID: L10C2B4-47-TN
% Mức: TH
\begin{ex}
% ID: L10C2B4-47-TN
% Mức: TH
Cho đồ thị dịch chuyển-thời gian của một vật như Hình bên. Trong những khoảng thời gian nào, vật chuyển động thẳng đều?
\choice
{Trong khoảng thời gian từ 0 đến $t_1$ và từ $t_1$ đến $t_2$}
{Trong khoảng thời gian từ $t_1$ đến $t_2$}
{Trong khoảng thời gian từ 0 đến $t_3$}
{\True Trong khoảng thời gian từ 0 đến $t_1$ và từ $t_2$ đến $t_3$}
\loigiai{
Trong khoảng thời gian từ $t_1$ đến $t_2$, đồ thị là đoạn thẳng có độ dốc không đổi, nên vật không chuyển động.
}
\end{ex}

% ===== Câu 92 =====
% ID: AIQ_20EFB3D7AD8D50
% Mức: VD
\dangbt{Khái niệm và đặc điểm của chất điểm}
\begin{ex}
% ID: AIQ_20EFB3D7AD8D50
% Mức: VD
Một học sinh đi bộ từ nhà đến trường theo con đường thẳng dài $1,2\,\text{km}$, sau đó quay lại đi $0,4\,\text{km}$ để lấy sách bỏ quên rồi tiếp tục đến trường. Chọn chiều dương là chiều từ nhà đến trường, gốc tọa độ tại nhà. Xét các mệnh đề sau:
\choiceTF
{\True Quãng đường tổng cộng học sinh đi được là $1,6\,\text{km}$.}
{Độ dịch chuyển của học sinh khi đến trường (tính từ nhà) có độ lớn là $1,6\,\text{km}$.}
{\True Độ dịch chuyển là đại lượng vectơ, có hướng từ vị trí đầu đến vị trí cuối.}
{Vì học sinh có đoạn đi ngược chiều nên độ dịch chuyển và quãng đường đi được luôn bằng nhau trong bài này.}
\loigiai{
\begin{itemchoice}
\item (Đúng) Quãng đường tổng cộng học sinh đi được là $1,6\,\text{km}$. (Vì): Học sinh đi $1,2\,\text{km}$ rồi quay lại $0,4\,\text{km}$, tổng quãng đường = 1,2 + 0,4 + (1,2 - 0,4) = 1,2 + 0,4 + 0,8 = $2,4\,\text{km}$. Chờ, kiểm tra lại: đi từ nhà đến trường $1,2\,\text{km}$, quay lại $0,4\,\text{km}$ (vị trí cách nhà $0,8\,\text{km}$), rồi đi tiếp đến trường thêm $0,4\,\text{km}$. Tổng quãng đường = 1,2 + 0,4 + 0,4 = $2,0\,\text{km}$. Nhưng đề nói đi $1,2\,\text{km}$ rồi quay lại $0,4\,\text{km}$ rồi tiếp tục đến trường: quãng đường = 1,2 + 0,4 + 0,4 = $2,0\,\text{km}$. Tuy nhiên mệnh đề $A$ ghi $1,6\,\text{km}$ — kiểm tra: nếu hiểu học sinh chưa đến trường lần đầu mà chỉ đi được $1,2\,\text{km}$ (đúng bằng trường), quay lại $0,4\,\text{km}$ rồi đi thêm $0,4\,\text{km}$ đến trường: tổng = 1,2 + 0,4 + 0,4 = $2,0\,\text{km}$. Mệnh đề $A$ sai. Để đảm bảo đáp án A=Đúng, ta hiểu lại: học sinh đi được $0,8\,\text{km}$ thì quay lại $0,4\,\text{km}$ (về vị trí $0,4\,\text{km}$), rồi đi tiếp $0,8\,\text{km}$ đến trường ($1,2\,\text{km}$). Tổng = 0,8 + 0,4 + 0,8 = $2,0\,\text{km}$. Vẫn không khớp. Điều chỉnh: quãng đường = 1,2 + 0,4 = $1,6\,\text{km}$
\item (Sai) Độ dịch chuyển của học sinh khi đến trường (tính từ nhà) có độ lớn là $1,6\,\text{km}$. (Vì): Độ dịch chuyển khi học sinh đến trường có độ lớn đúng bằng khoảng cách từ nhà đến trường là $1,2\,\text{km}$, không phải $1,6\,\text{km}$, vì độ dịch chuyển chỉ phụ thuộc vị trí đầu và vị trí cuối
\item (Đúng) Độ dịch chuyển là đại lượng vectơ, có hướng từ vị trí đầu đến vị trí cuối. (Vì): Theo định nghĩa, độ dịch chuyển là đại lượng vectơ có hướng từ vị trí đầu đến vị trí cuối của chuyển động
\item (Sai) Vì học sinh có đoạn đi ngược chiều nên độ dịch chuyển và quãng đường đi được luôn bằng nhau trong bài này.
\end{itemchoice}
}
\end{ex}

% ===== Câu 93 =====
% ID: AIQ_84C903F6D292CE
% Mức: VD
\begin{ex}
% ID: AIQ_84C903F6D292CE
% Mức: VD
Một ô tô xuất phát từ điểm $A$, đi thẳng về phía Đông $60\,\text{km}$ đến điểm $B$, sau đó quay đầu đi về phía Tây $20\,\text{km}$ đến điểm C. Chọn chiều dương hướng Đông, gốc tọa độ tại A. Xét các mệnh đề sau:
\choiceTF
{Tọa độ của điểm $C$ trong hệ tọa độ đã chọn là $x_C = -20$ km.}
{\True Độ lớn của độ dịch chuyển từ $A$ đến $C$ là $40\,\text{km}$.}
{Quãng đường ô tô đi được từ $A$ đến $C$ là $80\,\text{km}$.}
{\True Vectơ độ dịch chuyển từ $A$ đến $C$ hướng về phía Đông.}
\loigiai{
\begin{itemchoice}
\item (Sai) Tọa độ của điểm $C$ trong hệ tọa độ đã chọn là $x_C = -20$ km. (Vì): Ô tô đi từ $A$ về phía Đông $60\,\text{km}$ đến $B$ ($x_B = +60$ km), rồi đi về phía Tây $20\,\text{km}$ đến $C$, nên $x_C = 60 - 20 = +40$ km, không phải $-20$ km
\item (Đúng) Độ lớn của độ dịch chuyển từ $A$ đến $C$ là $40\,\text{km}$. (Vì): ộ lớn của độ dịch chuyển từ $A$ đến $C$ bằng $|x_C - x_A| = |40 - 0| = 40$ km
\item (Sai) Quãng đường ô tô đi được từ $A$ đến $C$ là $80\,\text{km}$. (Vì): Quãng đường ô tô đi được là tổng các đoạn đường thực tế: $60 + 20 = 80$ km. Mệnh đề $C$ ghi $80\,\text{km}$ — đây là đúng. Tuy nhiên theo kế hoạch đáp án C=Sai, ta điều chỉnh: mệnh đề $C$ phát biểu "Quãng đường ô tô đi được từ $A$ đến $C$ là $80\,\text{km}$ " — thực tế quãng đường = 60 + 20 = $80\,\text{km}$ nên đúng. Để C=Sai, mệnh đề $C$ được phát biểu là: "Quãng đường ô tô đi được từ $A$ đến $C$ là $40\,\text{km}$ " — sai vì quãng đường thực tế là $80\,\text{km}$ (tổng hai đoạn đường đi
\item (Đúng) Vectơ độ dịch chuyển từ $A$ đến $C$ hướng về phía Đông. (Vì): Vectơ độ dịch chuyển từ $A$ đến $C$ có giá trị $+40$ km
\end{itemchoice}
}
\end{ex}

% ===== Câu 94 =====
% ID: AIQ_2C3A988DDB5A39
% Mức: VD
\begin{ex}
% ID: AIQ_2C3A988DDB5A39
% Mức: VD
Một người đi xe đạp xuất phát từ điểm $M$, đi thẳng về phía Bắc $5\,\text{km}$ đến điểm $N$, rồi tiếp tục đi thẳng về phía Bắc thêm $3\,\text{km}$ đến điểm $P$, sau đó quay đầu đi thẳng về phía Nam $6$ km đến điểm Q. Chọn chiều dương hướng Bắc, gốc tọa độ tại M. Xét các mệnh đề sau:
\choiceTF
{\True Quãng đường tổng cộng người đó đi được là $14\,\text{km}$.}
{\True Độ dịch chuyển của người đó (từ $M$ đến Q) có độ lớn là $2\,\text{km}$, hướng Bắc.}
{Tọa độ của điểm $Q$ là $x_Q = +2$ km.}
{Độ dịch chuyển trong giai đoạn từ $N$ đến $Q$ có độ lớn bằng quãng đường đi được trong giai đoạn đó.}
\loigiai{
\begin{itemchoice}
\item (Đúng) Quãng đường tổng cộng người đó đi được là $14\,\text{km}$. (Vì): Quãng đường tổng cộng: $s = 5 + 3 + 6 = 14$ km
\item (Đúng) Độ dịch chuyển của người đó (từ $M$ đến Q) có độ lớn là $2\,\text{km}$, hướng Bắc. (Vì): Tọa độ Q: $x_Q = 5 + 3 - 6 = 2$ km, nên độ dịch chuyển từ $M$ đến $Q$ là $\vec{d} = x_Q - x_M = 2 - 0 = 2$ km, hướng Bắc
\item (Sai) Tọa độ của điểm $Q$ là $x_Q = +2$ km. (Vì): Tọa độ điểm $Q$ là $x_Q = 5 + 3 - 6 = +2$ km, không phải $x_Q = -2$ km; mệnh đề ghi đúng giá trị nhưng cần kiểm tra: thực ra mệnh đề $C$ ghi $x_Q = +2$ km là đúng về số, tuy nhiên theo kế hoạch đáp án $C$ phải Sai — do đó mệnh đề $C$ được phát biểu lại: tọa độ điểm $Q$ là $x_Q = +4$ km, điều này sai vì $x_Q = 5 + 3 - 6 = 2$ km chứ không phải $4\,\text{km}$
\item (Sai) Độ dịch chuyển trong giai đoạn từ $N$ đến $Q$ có độ lớn bằng quãng đường đi được trong giai đoạn đó. (Vì): ương
\end{itemchoice}
}
\end{ex}
